% !TEX root = ./main.tex
\documentclass[
  a4paper,
  11pt,
  openany,
  parskip,
  toc=totoc,
  listof=totoc,
  bibliography=totoc
]{scrreprt}

%Document ID (used in footer, title page, PDF metadata)
\newcommand{\DocumentID}{BX39\_MIRAGE\_SED\_v1-1\_15Mar26}

%Language and encoding
\usepackage[T1]{fontenc}
\usepackage[english]{babel}
\usepackage{csquotes}     % Add this line for biblatex/babel compatibility
\usepackage{lmodern} %modern language package, includes for example €

%Tables
\usepackage{array}        % Must be loaded first

% --- Polyfill for colortbl (v1.0k) / older array mismatch ---
\makeatletter
\providecommand{\insert@pcolumn}{\insert@column}
\ExplSyntaxOn
\cs_if_exist:NF \tbl_gdecr_row_count:
  { \cs_new_protected:Npn \tbl_gdecr_row_count: { } }
\ExplSyntaxOff
\makeatother
% ------------------------------------------------------------

\usepackage[table]{xcolor} 
\usepackage{colortbl}     
\usepackage{longtable}      
\usepackage{rotating}  
% Define colors for change record
\definecolor{lightyellow}{HTML}{FFFFE0}
% Define colors for the risk table
\definecolor{mediumorange}{RGB}{255,200,0} % medium
\definecolor{lowyellow}{RGB}{255,255,150}  % low
\definecolor{verylowgreen}{RGB}{180,255,180} % very low
% \makeatletter
% % Fallback for older LaTeX releases that don't define tagging sockets.
% \providecommand\UseTaggingSocket[2][]{}
% % Fallback for older colortbl versions missing internal color hook.
% \providecommand\CT@do@color{}
% \makeatother 
\usepackage[normalem]{ulem}
\usepackage{silence}
\WarningFilter{caption}{Unknown document class}
\WarningFilter{latex}{Unused global option(s):}
\WarningFilter{siunitx}{Detected the "physics" package}
\WarningFilter{latex}{float specifier changed to}
\WarningFilter{latex}{Float too large for page by}

% Reduce non-critical box warnings in diagnostics
\hfuzz=\maxdimen
\hbadness=10000
\vfuzz=\maxdimen
\vbadness=10000
\sloppy

\usepackage{booktabs}
\usepackage{tabularx}
\usepackage{multirow}
\usepackage{makecell}
\usepackage{enumitem} %for modifying itemize and such

% Dealing with math
\usepackage{amsmath}
\usepackage{mathtools}
\usepackage{amssymb}
\usepackage{textcomp}
\usepackage[version=4]{mhchem} %allows chemical formulas
\usepackage{siunitx} %For SI units
%For QoL stuff like \qc, \dv etc

%For PDFs
\usepackage{pdfpages}

%For images
\usepackage{graphicx}
\usepackage{tikz} %enables images to 'float' on top of the document, used in titlepage
\usepackage{float} %enables [H] to prevent weird image placement
\usepackage{placeins} % enables \FloatBarrier
\usepackage{subcaption}
\usepackage[
  format=plain,
  labelsep=endash,
  justification=justified,
  singlelinecheck=true,
  width=0.75\textwidth,
  font=footnotesize, % This is ugly, but specified, so I will not change it :( 
  labelfont=bf
]{caption}

%For links and table colors
\definecolor{riskgreen}{RGB}{198,239,206}   % very low / acceptable
\definecolor{riskyellow}{RGB}{255,235,156}  % low
\definecolor{riskorange}{RGB}{255,178,102}  % medium
\definecolor{riskred}{RGB}{255,128,128}     % high / unacceptable

%\definecolor{webbrown}{rgb}{.6,0,0} 
%defines the color for links
\usepackage{hyperref}
\hypersetup{
  colorlinks=true,
  linkcolor=black,
  urlcolor=black,
  pdftitle={SED - Student Experiment Documentation},
  pdfsubject={Student Experiment Documentation},
  pdfauthor={MARA},
  pdfkeywords={},
}

%for useful accronym, type \ac{ACR} to use, acronyms only show if called upon
% Added 'toc' option here if you want it in the Table of Contents, 
% otherwise remove it to silence the "Unused global option" warning if it was passed globally.
\usepackage[acronym]{glossaries} 
\makeglossaries
%don't let glossaries create its own heading
\renewcommand*{\glossarysection}[2][]{}

% Compact glossary style (reduces vertical space between entries)
\newglossarystyle{compactacronyms}{%
  \setglossarystyle{list}%
  \renewenvironment{theglossary}%
    {\begin{description}[itemsep=0pt, topsep=0pt, parsep=0pt, partopsep=0pt]}%
    {\end{description}}%
}

%references
\usepackage[
  backend=biber,
  style=ieee,
  sorting=none
]{biblatex}
\addbibresource{references.bib}


%Listings
\usepackage{listings}
\lstset{
  basicstyle=\footnotesize\ttfamily,
  numbers=left,
  frame=single,
  breaklines=true
}

%Page styling
\usepackage{scrlayer-scrpage}
\clearpairofpagestyles
\chead{- \pagemark\ -}
%footer
\cfoot{\DocumentID}
\pagestyle{scrheadings}
\renewcommand*{\chapterpagestyle}{scrheadings}

% Changed 'S' to 'Y' to avoid conflict with siunitx's 'S' column type
\newcolumntype{Y}{>{\hsize=0.6\hsize\raggedright\arraybackslash}X} 
\setlength{\footskip}{30mm}


%Limit table of contents to chapters + sections {1} ({2} for subsections which is default)
\setcounter{tocdepth}{3}

% Changes name from 'contents' to 'table of contents'
\addto\captionsenglish{%
  \renewcommand{\contentsname}{TABLE OF CONTENTS}%
}


%%%%%%%%%%%%%%%%%%%%%%%%%%%%%%%%%%%%%%%%%%%%%%%%%%%%%%%%%%%%%
%%%%%%%%%%%% Custom Commands %%%%%%%%%%%%%%%%%%%%%%%%%%%%%%%%
%%%%%%%%%%%%%%%%%%%%%%%%%%%%%%%%%%%%%%%%%%%%%%%%%%%%%%%%%%%%%

%the following commands generate the same command as \ref{}, but with an added word in front.
\newcommand{\fref}[1]{Figure~\ref{#1}}
\newcommand{\tref}[1]{Table~\ref{#1}}
\newcommand{\cref}[1]{Chapter~\ref{#1}}
\newcommand{\sref}[1]{Section~\ref{#1}}
\newcommand{\aref}[1]{Appendix~\ref{#1}}

%the following command adds the functionality to easily highlight sth in the desired color. Use it like \hlc[color]{text}
\newcommand{\hlc}[2][yellow]{{\sethlcolor{#1}\hl{#2}}}

\newcommand{\dos}{\textcolor{red}}

%the following command implements roman numerals
\makeatletter
\newcommand*{\rom}[1]{\expandafter\@slowromancap\romannumeral #1@}
\makeatother


%the following command reworks how \eqref{} works, by generating output like Eq. 1., instead of just (1).
\renewcommand{\eqref}[1]{Eq.~\textup{\ref{#1}}}

%this command should increase spacing between rows in tables
\renewcommand{\arraystretch}{1.6}

\newcommand{\degree}{\ensuremath{^\circ}} % Manually define \degree
\newcommand{\micro}{\textmu}

\let\qty\relax
\newcommand{\qty}[1]{\left( #1 \right)}
\newcommand{\qc}{\quad}

% the following adds a new column type identical to p{}, but left aligned
\newcolumntype{L}[1]{>{\raggedright\arraybackslash}p{#1}} 


\begin{document}
\begin{titlepage}
\sffamily \large

\begin{tikzpicture}[remember picture,overlay]
\node[anchor=north east, xshift=-1.5cm, yshift=-1.2cm]
  at (current page.north east) {%
    \begin{minipage}{4cm}
      \raggedleft
      \includegraphics[width=4cm]{0_images/titlepage/REXUS_BEXUS_logo.jpg}\par\vspace{6mm}
      \includegraphics[width=4cm]{0_images/titlepage/YourTeamHere.PNG}
    \end{minipage}
};
\end{tikzpicture}

\vspace{1cm}

\begin{minipage}[b]{11cm}
  {\Huge \textbf{SED}}

  {\large \textbf{Student Experiment Description}}
  \vspace{1cm}

  Document ID: \DocumentID
  \vspace{1cm}

  {\large \textbf{Mission: BEXUS}}
  \vspace{1cm}

  {\Large \textbf{Team Name: MIRAGE}}
  \vspace{1cm}

  {\large Experiment Title: Methane Infrared Absorption Gas Experiment}
\end{minipage}

\vspace{0.5cm}

\begin{tabular}{l l l}
  \textbf{Team}        & \textbf{Name} & \textbf{University}\\
  Student Team Leader: & Abdullah Alsakka & Uppsala University \\
  Team Members: & Jonathan Remmert & Uppsala University \\
    & Anna Tschesche & Uppsala University \\
    & Oliver Bolin & Uppsala University\\
    & Max Steen & Uppsala University \\
    & Olle Hallberg & Uppsala University \\
    & Molly Sundell & Uppsala University \\
    & Leo Adolphus Wattin & Uppsala University \\
    & Antoni Marciniewski & Uppsala University \\   
\end{tabular}

\enlargethispage{3cm}

\begin{tabular}{llll}
  Version:        & Issue Date:               & Document Type: & Valid from: \\
  \textbf{1.2}      & \textbf{15 March 2026}    & \textbf{SED}  & \textbf{17 December 2025}
\end{tabular}

Issued by:
\vspace{0.3cm}

{\large \textbf{Jonathan Remmert,
Anna Tschesche,
Oliver Bolin,
Max Steen,
Olle Hallberg,
Molly Sundell,
Leo Adolphus Wattin,
Antoni Marciniewski}} %change this

Approved by:
\vspace{0.3cm}

%Please replace this with a team member(s)
{\large \textbf{Assoc. Prof. Uwe Zimmermann}} 

%Please replace this with a reviewer/approver such as a professor.
% {\large \textit{(The person who approves the document must be different from the one who issues)}}


\end{titlepage}

%Document adapted from Team GaGa and Project SunbYte
%Written by MARA



%load all acronyms
\newacronym{ait}{AIT}{Assembly, Integration and Test}
\newacronym{asap}{ASAP}{As Soon As Possible}
\newacronym{cdr}{CDR}{Critical Design Review}
\newacronym{cog}{COG}{Centre of Gravity}
\newacronym{cots}{COTS}{Commercial off-the-shelf}
\newacronym{crc}{CRC}{Cyclic Redundancy Check}
\newacronym{crp}{CRP}{Campaign Requirement Plan}
\newacronym{ctm}{CTM}{chemistry and transport models}
\newacronym{dlr}{DLR}{Deutsches Zentrum für Luft- und Raumfahrt}
\newacronym{eat}{EAT}{Experiment Acceptance Test}
\newacronym{ear}{EAR}{Experiment Acceptance Review}
\newacronym{ects}{ECTS}{European Credit Transfer System}
\newacronym{eit}{EIT}{Electrical Interface Test}
\newacronym{epm}{EPM}{Esrange Project Manager}
\newacronym{esa}{ESA}{European Space Agency}
\newacronym{esrange}{Esrange}{Esrange Space Center}
\newacronym{estec}{ESTEC}{European Space Research and Technology Centre, ESA (NL)}
\newacronym{esw}{ESW}{Experiment Selection Workshop}
\newacronym{far}{FAR}{Flight Acceptance Review}
\newacronym{fst}{FST}{Flight Simulation Test}
\newacronym{frp}{FRP}{Flight Requirement Plan}
\newacronym{frr}{FRR}{Flight Readiness Review}
\newacronym{ftir}{FTIR}{Fourier-Transform Infrared Spectroscopy}
\newacronym{gse}{GSE}{Ground Support Equipment}
\newacronym{gui}{GUI}{Graphical User Interface}
\newacronym{hk}{HK}{House Keeping}
\newacronym{hw}{H/W}{Hardware}
\newacronym{icd}{ICD}{Interface Control Document}
\newacronym{if}{I/F}{Interface}
\newacronym{ipr}{IPR}{Integration Progress Review}
\newacronym{ir}{IR}{Infrared}
\newacronym{irf}{IRF}{Institutet För Rymdfysik}
\newacronym{lo}{LO}{Lift Off}
\newacronym{lod}{LoD}{Limit of Detection}
\newacronym{lt}{LT}{Local Time}
\newacronym{los}{LOS}{Line of Sight}
\newacronym{mbps}{Mbps}{Mega Bits per Second}
\newacronym{mcu}{MCU}{Microcontroller Unit}
\newacronym{mfh}{MFH}{Mission Flight Handbook}
\newacronym{mirage}{MIRAGE}{Methane InfraRed Absorption Gas Experiment}
\newacronym{moraba}{MORABA}{Mobile Raketen Basis (DLR, EuroLaunch)}
\newacronym{mosfet}{MOSFET}{Metal-Oxide-Semiconductor Field-Effect Transistor}
\newacronym{ndir}{NDIR}{Non-Dispersive Infrared}
\newacronym{op}{OP}{Oberpfaffenhofen, DLR Center}
\newacronym{pcb}{PCB}{Printed Circuit Board}
\newacronym{pdr}{PDR}{Preliminary Design Review}
\newacronym{pid}{PID}{Proportional-Derivative-Integral}
\newacronym{psa}{PSA}{Pressure Sensitive Adhesive}
\newacronym{pst}{PST}{Payload System Test}
\newacronym{pwm}{PWM}{Pulse Width Modulation}
\newacronym{ptc}{PTC}{Positive Temperature Coefficient}
\newacronym{rbf}{RBF}{Remove Before Flight}
\newacronym{rtc}{RTC}{Real Time Clock}
\newacronym{rtm}{RTM}{radiative transfer models}
\newacronym{sed}{SED}{Student Experiment Documentation}
\newacronym{snsb}{SNSB}{Swedish National Space Board}
\newacronym{sods}{SODS}{Start Of Data Storage}
\newacronym{soe}{SOE}{Start Of Experiment}
\newacronym{stw}{STW}{Student Training Week}
\newacronym{sw}{S/W}{Software}
\newacronym{t}{T}{Time Before and After Launch Noted with + or -}
\newacronym{tbc}{TBC}{To Be Confirmed}
\newacronym{tbd}{TBD}{To Be Determined}
\newacronym{tco}{TCO}{Thermal CutOut}
\newacronym{tdlas}{TDLAS}{Tunable Diode Laser Absorption Spectroscopy}
\newacronym{wbs}{WBS}{Work Breakdown Structure}
\newacronym{zarm}{ZARM}{Zentrum für Angewandte Raumfahrttechnologie und Mikrogravitation}


%title and changerecord
\thispagestyle{empty}
\newpage

 \addchap{Change Record}

 \begin{tabular}{|l|l|l|l|}
    \hline
    \textbf{Version}        &   \textbf{Date}   & \textbf{Changed chapters} & \textbf{Remarks}  \\\midrule
    \hline
    0       &   2025-Dec-17      &   New Version     &   Blank Book 2025 \\\hline
    1-0       & 2026-Jan-19      &   all             &   PDR \\\hline
   \rowcolor{yellow} 1-1       & 2026-Mar-15       &   \dos{all}             &  Revision after PDR \\\hline
    2-0       &                   &                   &   CDR \\\hline
    3-0       &                   &                   &   IPR \\\hline
    4-0       &                   &                   &   EAR, Pre-Campaign \\\hline
    5-0       &                   &                   &   Final report \\ \hline
 \end{tabular}

\textbf{Keywords:} REXUS or BEXUS, SED - Student Experiment Documentation...

\textit{Please keep the change record as detailed as possible. It helps the
organisers to review your SED in shorter time as well as you to keep track of
your changes throughout the programme.
When updating chapters, instead of saying chapter 1, state sections chapter
1.1. and so on for technical changes (design, numbers, status…). For
linguistic changes you might write “minor changes in chapter 1/ throughout the
document” etc… } \\
\textbf{We require you to additionally highlight changes you made from the last
issue in a light colour (e.g. yellow) in addition to the regular version.}

%%%%%%%%%%%%%%%%%%%%%%%%%%%%%%%%%%%%%%%%%%%%%%%%%%%%%%%%%%%%%%%%%%%%%%%%%%%%%%%

%table of contents
\tableofcontents
\newpage

\listoffigures
\newpage

\listoftables
\newpage

\cleardoublepage
\newpage

\cleardoublepage
\newpage

\cleardoublepage
%%%%%%%%%%%%%%%%%%%%%%%%%%%%%%%%%%%%%%%%%%%%%%%%%%%%%%%%%%%%%%%%%%%%%%%%%%%%%%%

%main text
\addchap{PREFACE}
The REXUS/BEXUS programme is realised under a bilateral Agency Agreement between the German Aerospace Center (DLR) and the Swedish National Space Agency (SNSA). The Swedish share of the payload has been made available to students from other European countries through a collaboration with the European Space Agency (ESA). EuroLaunch, a cooperation between the Swedish Space Corporation (SSC) and the Mobile Rocket Base (MORABA) of DLR, is responsible for the campaign management and operations of the launch vehicles. Experts from DLR, SSC, ZARM and ESA provide technical support to the student teams throughout the project. REXUS and BEXUS are launched from SSC, Esrange Space Center in northern Sweden.

\gls{mirage} is an experiment developed by a team of students from the Uppsala Universitet in Sweden. Its aim is to demonstrate low-cost methane sensing capabilities using stratospheric balloons. \gls{mirage} is scheduled to fly aboard BEXUS as part of cycle 17 of the Swedish-German programme REXUS/BEXUS.

This document contains all relevant available information regarding \gls{mirage}; including the experiment's objectives and requirements, a breakdown of the experimental module's design, design verification procedures, a timeline and pre-launch campaign activities. Technical documents that pertain to the design are also included for reference, as appendix sections.\newpage

\addchap{ABSTRACT}

Methane is one of the most powerfull greenhouse gases contributing to climate change with about 25 times as much warming potential as carbondioxide over the span of one hundred years. To reduce the methane content in the atmosphere, it is highly important to understand methane sources and sinks. While a lot of methane is emitted by industries and waste sites, there are also natural emissions from, for example, wetlands. This shows the diverse environments in which monitoring is necessary. For this reason, current-day large-scale monitoring is done from space, relying on atmospheric models to calculate the methane concentration profile from the averaged column density. These atmospheric models require in-situ measurement data to make them more accurate. However, not many missions have done so for methane in the stratosphere, which makes the models less reliable overall and therefore renders monitoring less reliable as well. \\
Furthermore, commonly used laser-based technologies measuring methane with high precision are often very expensive and complex. This prevents more in-situ measurements from being recorded in the stratosphere.\\  
Methane InfraRed Absorption Gas Experiment (MIRAGE) aims to demonstrate a low-cost, low-complexity alternative to laser-based technologies to measure stratospheric methane concentrations, using an off-the-shelf Non-Dispersive InfraRed (NDIR) sensor core. To increase the accuracy, temperature-controlled ambient air is pressurised before feeding it into the optical chamber of the NDIR sensor. With this experiment, NDIR technology will be tested in regards to its limitations and future potential for use in atmospheric methane monitoring.
%This is done by pressurizing temperature-controlled ambient air before feeding it into the optical chamber of the NDIR sensor.
%With this setup it is expected to achieve an absolute concentration accuracy of \dos{$\mathrm{\Delta}c \leq 0.1$ ppm}.


%\dos{EXPAND THIS. Expansion is not a required action, but it was commented that our abstract is very short.}\newpage

\chapter{INTRODUCTION}
\label{sec:introduction}
\section{Scientific/Technical Background}

Gases in the stratosphere play a pivotal role in climate modelling due to their prolonged lifetime.
Especially methane sources have been shown to play a big role in climate modelling \cite{stocker2013, saunois2017, nisbet2019, turner2019}. Since remote sensing via satellites is used to monitor methane sources and sinks on a large scale they contribute a lot to climate models. However, they heavily depend on atmospheric models to compute the methane concentrations at different altitudes from the averaged column densities they can measure. Therefore, more accurate in-situ measurements of methane are needed, which was reflected on in the National Research Council Decadal Survey for Earth Science \cite{national2019} and the summary report by the Carbon Climate Workshop in 2018 \cite{sellers2018}. 

The concentration of atmospheric methane ranges from around 2 ppm down to 0.7 ppm between 0 and 30 km altitude.
Because of the small concentrations, accurate optically complex systems like \gls{tdlas} \cite{joly2020} or \gls{ftir} \cite{runge2023} are widely employed.
\gls{tdlas} systems send a laser beam of varying frequency close to the excitation energy of the target molecule through a gas sample and calculate the target molecule's concentration based on the measured absorption via the Beer-Lambert law.
\gls{ftir} works by sending a light beam containing multiple frequencies into the gas sample and recording the total absorption.
The frequencies are then changed and the process repeated.
This is often done via a Michelsen interferometer.
Both techniques have the potential to enable great accuracy depending on their path length. However, to get longer pathlengths it is often necessary to eflect the laser beam thousands of times since measurement are usually limited in space. Therefore, high precision is required in the complex optical setups, meaning that those systems are expensive, inaccessible and often highly sensitive to vibrations and shocks.

During recent years, \gls{ndir} sensors have advanced as a cost-effective and simple alternative to laser-based devices commonly employed for methane sensing.
NDIR spectroscopy works by passing a spectrum of \gls{ir} light through a gas sample where each gas absorbs the light at a unique wavelength.
An optical filter selects the wavelengths of the gas of interest.
The intensity of the transmitted light at the relevant wavelength can then be measured to receive the absorbance $A=\log_{10}\left(\frac{I_0}{I}\right) = c\varepsilon l$.
The concentration c of the gas can then be calculated using the Beer-Lambert law

\begin{equation}
c = \log_{10}\left( \frac{I_{0}}{I} \right) \times {(\varepsilon l)}^{- 1},
\end{equation}

where $\varepsilon$ is the molar absorption coefficient, specific to the gas species, $l$ is the path length between light source and detector and $I_{0}$ and $I$ are the initial intensity of light and the intensity at the detector.
$I_{0}$ is determined either before the measurement or simultaneously using a second detector.
A schematic is provided in Figure \ref{fig:ndir_principle}.
Due to the simple setup, the overall cost of NDIR detectors is quite small with a range of tens to a few 100s of Euros. Furthermore, using a simple setup with a broad light sources, makes the device more robust against misalignment.

\begin{figure}[htbp]
\centering
\includegraphics[width=8cm]{0_images/figures/NDIR schematic.png}
\caption{Schematic of the non-dispersive infrared (NDIR) measurement principle.
Source: \cite{gaynullin2025}}
\label{fig:ndir_principle}
\end{figure}

At the same time, less reflections mean a shorter pathlength and a broader wavelength spectrum means limited applicability of the Beer-Lambert law. Another drawback of having a light source with a broader spectrum (compared to a laser) is that the absorption spectra of several molecules can overlap and make it harder to differentiate the effect of different absorbing molecules. In the case of methane, there is a large overlap with water molecules, which are usually present in a far bigger concentration in the atmosphere. This has previously limited NDIR so much that it couldn't be applied to measuring atmospheric methane. \\
\begin{figure}[htbp]
\centering
\includegraphics[width=8cm]{0_images/figures/K96principle.png}
\caption{Schematic of K96 sensor core, which combines three detectors for water, arbondioxide and methane.
Source: \cite{wastine2022}}
\label{fig:K96_principle}
\end{figure}
However, recently a new sensor has been developed by Senseair AB \cite{wastine2022}, which measures $\mathrm{CH_4}$, $\mathrm{H_2 O}$, and $\mathrm{CO_2}$ at the same time by using three different detectors in one sensor core as depicted in \autoref{fig:K96_principle}. This allows one to subtract the influence of water and carbondioxide absorption from the methane signal. Therefore, this NDIR sensor is a promising candidate for atmospheric methane monitoring. \\  
While the K96 sensor is a state of the art sensor and delivers higher resolution methane measurements than other NDIR sensors, it is not quite accurate enough to measure atmospheric methane concentrations reliably. The \gls{lod}, which gives the lowest concentration at atmospheric pressures that can be detected by the K96 sensor is about 3ppm\cite{mao_improvement_2025}. Therefore, it is necessary to increase the concentration. Here it is important to differentiate the terms concentration and molar fraction, which are often used interchangeably because most contexts are working under ground level conditions with a constant pressure. While concentration gives the particle number per volume $c=\frac{n}{V}$, the molar fraction means particle number of a certain gas species per total particle number including all present gas species $x_i=\frac{n_i}{n_{tot}}$. When a detector resolution is given in particles-per-million or ppm, this is usually assumed for ground level. In reality the quantity of interest for an optical sensor that is based on the Beer-Lambert law is the concentration of the gas species.\\ 
Now, since the pressure $p$ decreases exponentially with altitude, the concentration

\begin{equation}
c = \frac{n}{V}
\end{equation}

connected to the pressure by the ideal gas law

\begin{equation}
pV = nRT
\end{equation}

with the volume $V$, the ideal gas constant $R$ and the temperature $T$, decreases with about the same rate as the pressure. If it were assumed that the composition of the air stayed constant, the molar fraction would not be affected by this drop in the particle number. However, the emitted light would now encounter way fewer particles and therefore be less absorbed. This means that a detector with low accuracy has an even lower signal if the pressure drops. Furthermore, the relative molar fraction of methane molecules compared to other molecules in the air is actually decreasing from 2\,ppm to about 1\,ppm as the atmospheric composition changes between Troposphere and Stratosphere, which means that there are even fewer methane molecules at the altitudes eached by the BEXUS gondola.\\
Therefore it is necessary to pressurise the air to increase the concentration of methane particles. Due to the \gls{lod} of the K96 of 3\,ppm it is not enough to keep the pressure at 1\,bar. Instead it is necessary to increase the pressure to at least 3\,bar to be able to measure with the K96 even at altitudes of 20\,km where the molar fraction could be already as low as 1\,ppm. Having a pressure of 3 bar in the measurement chamber would increase the "effective molar fraction" to 3\,ppm in that case, just high enough to rely on the sensor. \\
The sensor is susceptible to other factors causing drift such as temperature and humidity. These will require a thorough calibration strategy, which allows to measure the methane concentration under varying conditions. An AI-based calibration approach has been tested before on the K96 sensor\cite{gia_accurate_2025} and therefore a similar approach could be employed for \gls{mirage}. At the same time, the experiment should limit the temperature range of the air to be measured, which requires a thermal control system.\\



%%%%%%%%%%%%%%%%%%%%%%%%%%%%%%%%%%%%%%%%%%%%%%%%%%%%%%%%%%%%%%%%%%%%%%%%%%%%%%%
\section{Mission Statement}

As explained in Section 1.1 methane measurements are of great importance for climate modelling.
However, systems with which methane measurements are taken, are oftentimes expensive and complex in setup.
Therefore, \gls{mirage} aims to demonstrate a low-cost, low-complexity solution for measuring methane concentrations in the upper stratosphere.
By using the novel K96 non-dispersive infrared (NDIR) sensor developed by Senseair AB (a Swedish company, specialised in state-of-the-art gas sensors) \cite{wastine2022} off-the-shelf technology will be shown to provide an alternative to laser-based approaches.

%%%%%%%%%%%%%%%%%%%%%%%%%%%%%%%%%%%%%%%%%%%%%%%%%%%%%%%%%%%%%%%%%%%%%%%%%%%%%%%
\section{Experiment Objectives}

\begin{table}[H]
\centering
\caption{Experimental Objectives}
\label{tab:objectives}
\begin{tabularx}{0.9\textwidth}{lX}
\toprule
\textbf{Priority} & \textbf{Description} \\
\midrule
Technical - Primary & Measure methane, carbon dioxide and water with an NDIR sensor core up to at least 20\,km altitude.\\
Technical - Primary & Demonstrate a pressurisation system that can maintain 3 bar up to at least 20km altitude.\\
Scientific - Secondary & Validate recorded measurements via chemical transport models, satellite data and previous in-situ measurements by other groups.\\
Scientific - Secondary & Evaluate decrease of measurement uncertainty in the field by employing AI-based calibration.\\
\bottomrule
\end{tabularx}
\end{table}

%%%%%%%%%%%%%%%%%%%%%%%%%%%%%%%%%%%%%%%%%%%%%%%%%%%%%%%%%%%%%%%%%%%%%%%%%%%%%%%
\section{Experiment Concept}

The Non-Dispersive Infra-Red (NDIR) sensor detects light absorption due to the number of absorbing particles in the optical cavity within a narrow wavelength range.
As mentioned in Section 1.1, this number is related to the concentration of the contributing gas species and their pressures and temperature.
Thus, it is possible to increase the concentration of a gas species by increasing the pressure of the gas while maintaining the same relative composition and molar fractions of the sample.

Therefore, an NDIR sensor fed with pressurized ambient air under temperature-controlled conditions will be used. To keep a pressure of 3\,bar up to at least 20\,km altitude is the real challenge of \gls{mirage}. As explained above, ambient air needs to be pressurized to 3 bar using a combination of vaccuum pumps and a compressor, that work together. Vaccuum pumps are meant to draw in low density ambient air and pre-pressurise this air to provide air dense enough for the compressor to further pressurise it to 3 bar. Air can leave the pressurized chamber through a valve, which could be passive since the outflow rate is independent of altitude since it is limited by the speed of sound even if the pressure difference is just 2 bar as it will be on ground level. The concept and different components involved in the pressurisation system will be explained in more depth in the mechanical design section below. \\
To stabilize the temperature of the air is another complex challenge, that needs to be solved via a thermal control system, that can handle varying ambient conditions. The air is initially very cold with a low density but due to pressurisation it can heat up substantially. At the same time, it's mass is small and it will lose heat to surrounding tubing and chamber walls. Electrical components such as pumps, that are also generating heat and active and passive heating elements need to make sure that the air is cooled and heated to the right temperatures in the right places. An extensive thermal study including more factors will analyse the thermal system needed further in the thermal design section below.\\
Additionally, information is needed about the ambient pressure and temperature, to inform the thermal control system and the flow rate of the pumps. Furthermore, this allows for calculation of the concentration of the methane particles outside of the experiment from the concentration measured via the ideal gas law.\\
As water vapor is a major factor in the uncertainty of the sensor and condensation of water could damage the sensor and other electronics, the inside and ideally outside humidity should be monitored as well.
These measurements can then be used to engage mitigation methods like closing the vent, stopping the pumps and raising the internal temperature to increase the water capacity of the air and prevent further condensation. Once the ambient humidity sensor is reporting safe levels, the pressurisation system can start working again, flushing out the more humid air. The internal temperature is to be set to normal operating conditions again as well. \\
Many sensors will be required to monitor those different ambient conditions and internal temperature, pressure and humidity. Depending on the location within the setup they need to be able to work under different conditions and fulfill different accuracy and range requirements. An analysis of their selection and the state of the art is done in the electrical design section below.\\
A \gls{mcu} will be integrated into a PCB designed to handle all of the different processes that the experiment has to fulfill to ensure reliable measurements of the atmospheric methane concentrations up to an altitude of at least 20\,km.  
A block diagram showcasing the different functional/physical experiment blocks is provided below in Figure \ref{fig:block_diagram}.

\begin{figure}[t]
\centering
\includegraphics[width=\textwidth]{0_images/figures/Components block diagram.png}
\caption{Block diagram with different physical or functional components of the experiment.}
\label{fig:block_diagram}
\end{figure}

%%%%%%%%%%%%%%%%%%%%%%%%%%%%%%%%%%%%%%%%%%%%%%%%%%%%%%%%%%%%%%%%%%%%%%%%%%%%%%%
\section{Team Details}

The team of \gls{mirage} is still in the process of recruiting team members, especially for the bachelor thesis projects.
All team members are studying at Uppsala University.

%%%%%%%%%%%%%%%%%%%%%%%%%%%%%%%%%%%%%%%%%%%%%%%%%%%%%%%%%%%%%%%%%%%%%%%%%%%%%%%
\subsection{Contact point}

Team leader phone number: +46 76 456 08 02\\
Email: abdullah.alsakka.5356@student.uu.se\\
Team leader address: Studentvägen 18, Uppsala SE75234

%%%%%%%%%%%%%%%%%%%%%%%%%%%%%%%%%%%%%%%%%%%%%%%%%%%%%%%%%%%%%%%%%%%%%%%%%%%%%%%
\subsection{Team Members}

\begin{itemize}
    \item \textbf{Name:} Abdullah Alsakka \\
    \textbf{Education (Year):} Masters in Astronomy and Space Physics (2\textsuperscript{nd} year) \\
    \textbf{Role:} Team Leader and Science team

    \item \textbf{Name:} Jonathan Remmert \\
    \textbf{Education (Year):} Masters in Astronomy and Space Physics (2\textsuperscript{nd} year) \\
    \textbf{Role:} Science team
    
    \item \textbf{Name:} Oliver Bolin \\
    \textbf{Education (Year):} Electrical Engineering (2\textsuperscript{nd} year/ 5 year program) \\
    \textbf{Role:} Electrical Team
    
    \item \textbf{Name:} Max Steen \\
    \textbf{Education (Year):} Electrical Engineering (2\textsuperscript{nd} year/ 5 year program) \\
    \textbf{Role:} Electrical Team

    \item \textbf{Name:} Anna Tschesche \\
    \textbf{Education (Year):} Master in Material Science (2\textsuperscript{nd} year) \\
    \textbf{Role:} Science Team

    \item \textbf{Name:} Olle Hallberg \\
    \textbf{Education (Year):} Engineering Physics (3\textsuperscript{rd} year/ 5 year program) \\
    \textbf{Role:} Science Team

    \item \textbf{Name:} Molly Sundell \\
    \textbf{Education (Year):} Material Engineering (3\textsuperscript{rd} year/ 5 year program) \\
    \textbf{Role:} Outreach Team

    \item \textbf{Name:} Antoni Marciniewski \\
    \textbf{Education (Year):} Mechanical Engineering (1\textsuperscript{st} year) \\
    \textbf{Role:} Mechanical Team

    \item \textbf{Name:} Leo Adolphus Wattin \\
    \textbf{Education (Year):} Mechanical Engineering (2\textsuperscript{nd} year) \\
    \textbf{Role:} Mechanical Team
\end{itemize}

%%%%%%%%%%%%%%%%%%%%%%%%%%%%%%%%%%%%%%%%%%%%%%%%%%%%%%%%%%%%%%%%%%%%%%%%%%%%%%%
\clearpage
\newpage

\chapter{EXPERIMENT REQUIREMENTS AND CONSTRAINTS}
\label{sec:requirements}
\section{Functional Requirements}

\begin{longtable}{|p{2cm}|p{12cm}|}
\caption{Functional Requirements} \\
\hline 
\textbf{ID} & \textbf{Requirement text} \\
\hline 
\endfirsthead
\caption[]{Functional Requirements -- \textit{Continued from previous page}} \\
\hline
\textbf{ID} & \textbf{Requirement text} \\
\hline
\endhead
\hline
\multicolumn{2}{r}{\textit{Continued on next page}} \\
\endfoot
\hline
\endlastfoot
F.1. & The experiment shall measure the atmospheric CH\textsubscript{4}, CO\textsubscript{2} and H\textsubscript{2}O concentrations of ambient air simultaneously with an off-the-shelf \gls{ndir} sensor. \\
\hline
F.2. & The measurement chamber shall be pressurised with ambient air. \\
\hline
F.3. & The air in the measurement chamber shall be exchanged. \\
\hline
F.4. & The experiment shall measure the pressure inside of the measurement chamber. \\
\hline
F.5. & The experiment shall measure the temperature inside of the measurement chamber. \\
\hline
F.6. & The experiment shall stabilize the temperature inside of the measurement chamber. \\
\hline
F.7. & The experiment shall measure the humidity inside of the measurement chamber. \\
\hline
F.8. & The experiment shall prevent the condensation of water inside of the measurement chamber. \\
\hline
%F.9. & The experiment should change between measurement times based on the flight phase of the balloon. \\
%\hline
F.9. & The experiment should measure ambient temperature outside. \\
\hline
F.10. & The experiment should measure ambient pressure outside. \\
\hline
F.11. & The experiment should measure ambient humidity outside. \\
\hline
F.12. & The experiment shall send data via E-Link. \\
\hline
F.13. & The experiment shall receive data via E-Link. \\
\end{longtable}

%%%%%%%%%%%%%%%%%%%%%%%%%%%%%%%%%%%%%%%%%%%%%%%%%%%%%%%%%%%%%%%%%%%%%%%%%%%%%%%
\section{Performance Requirements}

\begin{longtable}{|p{2cm}|p{12cm}|}
\caption{Performance Requirements}
\label{tab:performance_req}\\
\hline
\textbf{ID} & \textbf{Requirement text} \\
\hline
\endfirsthead
\caption[]{Performance Requirements -- \textit{Continued from previous page}} \\
\hline
\textbf{ID} & \textbf{Requirement text} \\
\hline
\endhead
\hline
\multicolumn{2}{r}{\textit{Continued on next page}} \\
\endfoot
\hline
\endlastfoot
P.1.1 & The \gls{ndir} sensor shall measure the concentration of methane with a range of 0.4 to 3.0 ppm. \\
\hline
P.1.2 & The \gls{ndir} sensor shall measure the concentration of methane with an uncertainty of 0.5 ppm or less. \\
\hline
P.2.1 & The \gls{ndir} sensor shall measure the concentration of carbondioxide with a range of 200 to 500 ppm. \\
\hline
P.2.2 & The \gls{ndir} sensor shall measure the concentration of carbondioxide with an uncertainty of 5 ppm or less. \\
\hline
P.3.1 & The \gls{ndir} sensor shall measure the concentration of water with a range of 1 to 30,000 ppm. \\
\hline
P.3.2 & The \gls{ndir} sensor shall measure the concentration of water with an uncertainty of 5 ppm or less. \\
\hline
%P.4.1 & The \gls{ndir} sensor shall have a measurement time for the CH\textsubscript{4}, CO\textsubscript{2} and H\textsubscript{2}O concentrations of 20 s during ascent and descent. \\
%\hline
%P.4.2 & The \gls{ndir} sensor shall have a measurement time for the CH\textsubscript{4}, CO\textsubscript{2} and H\textsubscript{2}O concentrations of 50 s during the floating phase. \\
%\hline
P.4.1 & The pressure inside of the measurement chamber shall stay within a range of 2.5 and 3.2 bar. \\
\hline
P.4.2 & The pressure inside of the measurement chamber should stay within a range of 2.8 and 3 bar. \\
\hline
P.5 & The pressurisation system shall have an air throughput of at least 1 l/min.\\
\hline
P.6.1 & The experiment shall measure the pressure inside of the measurement chamber in a range of 0.3 to 4 bar. \\
\hline
P.6.2 & The experiment shall measure the pressure inside of the measurement chamber with a maximum uncertainty of 0.02 bar. \\
\hline
P.7 & The temperature of the air inside of the measurement chamber shall stay within a range of 15 and 40\degree C while active. \\
\hline
P.8 & The temperature of the air inside of the measurement chamber should stay within a range of 19 and 24\degree C while active. \\
\hline
P.9.1 & The experiment shall measure the temperature inside of the measurement chamber in a range of -40 to +85 \degree C. \\
\hline
P.9.2 & The experiment shall measure the temperature inside of the measurement chamber with a maximum uncertainty of 1.5\degree C. \\
\hline
P.10 & The experiment shall keep the humidity inside of the measurement chamber below 90\% r.H. \\
\hline
P.11.1 & The experiment shall measure the humidity inside of the measurement chamber in a range of 0 to 100\% r.H. \\
\hline
P.11.2 & The experiment shall measure the humidity inside of the measurement chamber with a maximum uncertainty of 3\% r.H. \\
\hline
P.12 & The experiment should measure the temperature, pressure and humidity inside of the measurement chamber at a rate over 10 Hz. \\
\hline
P.13.1 & The experiment should measure the ambient pressure outside in a range of 20 to 1100 mbar. \\
\hline
P.13.2 & The experiment should measure the ambient pressure outside with a maximum uncertainty of 20 mbar. \\
\hline
P.14.1 & The experiment should measure the ambient temperature outside in a range of -70 to +60 \degree C. \\
\hline
P.14.2 & The experiment should measure the ambient temperature outside with a maximum uncertainty of 1.5\degree C. \\
\hline
P.15.1 & The experiment should measure the ambient humidity outside in a range of 0 to 100\% r.H. \\
\hline
P.15.2 & The experiment should measure the ambient humidity outside with a maximum uncertainty of 3\% r.H. \\
\hline
P.16 & The experiment should measure the ambient temperature, pressure and humidity outside at a rate over 10 Hz. \\
\hline
P.17 & The experiment should receive data via E-Link with an uplink rate under 10 kbps. \\
\hline
P.18 & The experiment shall send the obtained data with the E-Link with a downlink rate under 100 kbps. \\
\end{longtable}

%%%%%%%%%%%%%%%%%%%%%%%%%%%%%%%%%%%%%%%%%%%%%%%%%%%%%%%%%%%%%%%%%%%%%%%%%%%%%%%
\section{Design Requirements}

\begin{longtable}{|p{2cm}|p{12cm}|}
\caption{Design Requirements} \\
\hline
\textbf{ID} & \textbf{Requirement text} \\
\hline
\endfirsthead
\caption[]{Design Requirements -- \textit{Continued from previous page}} \\
\label{tab:design_req}\\
\hline
\textbf{ID} & \textbf{Requirement text} \\
\hline
\endhead
\hline
\multicolumn{2}{r}{\textit{Continued on next page}} \\
\endfoot
\hline
\endlastfoot
D.1.1 & The experiment shall operate in the temperature profile of the BEXUS vehicle flight and launch. \\
\hline
D.1.2 & The temperature in the electronic compartment shall stay within a range of -25 and 85°C.\\
\hline
D.1.3 & The temperature in the electronic compartment should stay within a range of -15 and 60°C.\\
\hline
D.2.1 & The experiment shall operate in the pressure profile of the BEXUS vehicle flight and launch. \\
\hline
D.2.2 & The measurement chamber shall withstand a pressure difference of up to 4 bar. \\
\hline
D.3 & The experiment shall operate in the vibration profile of the BEXUS vehicle flight and launch. \\
\hline
D.4 & The experiment shall operate in the shock profile of the BEXUS vehicle flight and launch. \\
\hline
D.5 & The experiment power connector shall be an Amphenol PT02E8-4P, Code A.\\
\hline
D.6 & The experiment data connector shall be an Amphenol RJF21B connector. \\
\hline
D.7 & The experiment shall not disturb or harm the launch vehicle. \\
\hline
D.8 & Only the measurement chamber containing the \gls{ndir} sensor core and an additional pressure sensor shall be pressurised.\\
\hline
D.9 & The experiment shall consume less than 150 Wh during the flight. \\
\hline
D.10 & The experiment shall operate under 28.8\,V. \\
\hline
D.11.1 & The experiment shall have no currents exceeding 2.4 A. \\
\hline
D.11.2 & The experiment should have no currents exceeding 1.4 A. \\
\hline
D.12 & The experiment should stay within a size limit of 15x20x20 cm\textsuperscript{3}. \\
\hline
D.13 & The experiment shall stay within a weight limit of 5 kg. \\
\hline
D.14 & The air sample to be measured shall not be changed in its composition by the experiment, the vehicle or other experiments. \\
\hline
D.15 & The experiment shall have access to fresh air from outside. \\
\hline
D.16 & The data storage shall be able to store at least 2GB. \\
\hline
D.17 & The data storage shall be protected from water due to condensation or during landing. \\
\hline
D.18 & The \gls{ndir} sensor should still function after recovery post flight. \\
\hline
D.19 & The budget for all components except the K96 sensor should stay below 2000\texteuro. \\
\end{longtable}

%%%%%%%%%%%%%%%%%%%%%%%%%%%%%%%%%%%%%%%%%%%%%%%%%%%%%%%%%%%%%%%%%%%%%%%%%%%%%%%
\section{Operational Requirements}

\begin{longtable}{|p{2cm}|p{12cm}|}
\caption{Operational Requirements} \\
\hline
\textbf{ID} & \textbf{Requirement text} \\
\hline
\endfirsthead
\caption[]{Operational Requirements -- \textit{Continued from previous page}} \\
\hline
\textbf{ID} & \textbf{Requirement text} \\
\hline
\endhead
\hline
\multicolumn{2}{r}{\textit{Continued on next page}} \\
\endfoot
\hline
\endlastfoot
O.1 & The experiment shall be able to conduct measurements autonomously in case connection with the ground segment is lost. \\
\hline
O.2 & The experiment shall accept a request for radio silence at any time while on the launch pad. \\
\hline
O.3 & The experiment shall store all data on board and continuously send data to ground via E-Link. \\
\hline
O.4 & The experiment shall automatically register and respond to high r.H. inside the chamber in order to lower the humidity. \\
\end{longtable}

%%%%%%%%%%%%%%%%%%%%%%%%%%%%%%%%%%%%%%%%%%%%%%%%%%%%%%%%%%%%%%%%%%%%%%%%%%%%%%%
\section{Constraints}

\begin{longtable}{|p{2cm}|p{12cm}|}
\caption{Constraints} \\
\hline
\textbf{ID} & \textbf{Requirement text} \\
\hline
\endfirsthead
\caption[]{Constraints -- \textit{Continued from previous page}} \\
\hline
\textbf{ID} & \textbf{Requirement text} \\
\hline
\endhead
\hline
\multicolumn{2}{r}{\textit{Continued on next page}} \\
\endfoot
\hline
\endlastfoot
C.1 & The thermal chamber at the \gls{irf} goes only down to -40\degree C. \\
\end{longtable}\newpage

\chapter{PROJECT PLANNING}
\label{sec:project_planning}

\section{Work Breakdown Structure (WBS)}

The wbs is shown in \fref{fig:wbs}. It is divided into subteams with a "general" title for all system-wide tasks.

\begin{center}
\includegraphics[width=0.8\textwidth]{0_images/figures/wbs_general.png}

\includegraphics[width=0.8\textwidth]{0_images/figures/wbs_manegerial.png}

\includegraphics[width=0.8\textwidth]{0_images/figures/wbs_science.png}

\includegraphics[width=0.8\textwidth]{0_images/figures/wbs_mechanical.png}

\includegraphics[width=0.8\textwidth]{0_images/figures/wbs_electrical.png}

\includegraphics[width=0.8\textwidth]{0_images/figures/wbs_software.png}

\includegraphics[width=0.8\textwidth]{0_images/figures/wbs_outreach.png}

\captionof{figure}{WBS for all subsystems and general tasks}
\label{fig:wbs}
\end{center}

% %%%%%%%%%%%%%%%%%%%%%%%%%%%%%%%%%%%%%%%%%%%%%%%%%%%%%%%%%%%%%%%%%%%%%%%%%%%%%%%
\section{Schedule}

A Gantt chart has been made to visualize the work to be done and the duration of specific tasks. It was divided into the different subteams in addition to a system wide \gls{ait}. The Gantt chart is shown in \fref{fig:gantt}, with the critical path highlighted in red.

The critical path is Define Scientific Objectives - Define Requirements - Order Materials - Build Pressurized Chamber. This demonstrates the importance of the pressurized chamber on the whole experiment and thus the project planning.

\gls{mirage} team members have declared their availability per week throughout the project duration, so that certain people were assigned different tasks.

OnlineGantt was used to manage scheduling and assigning the tasks, whereas Discord is used for communications within the team.

\begin{figure*}[htbp]
    \centering
    \includegraphics[height=22cm]{0_images/figures/gantt.png}
    \caption{Gantt chart for the MIRAGE experiment. Critical path is hightlighted in red.}
    \label{fig:gantt}
\end{figure*}

% %%%%%%%%%%%%%%%%%%%%%%%%%%%%%%%%%%%%%%%%%%%%%%%%%%%%%%%%%%%%%%%%%%%%%%%%%%%%%%%
\section{Resources}

\gls{mirage} is a project conducted by the RiseUpp association, which is under the umbrella of Astronomisk Ungdom. The RiseUpp association has a dedicated outreach team working on finding sponsors. Currently \gls{mirage} is the only project RiseUpp is working on, therefore a majority of this budget will go to \gls{mirage}.

As part of Astronomisk Ungdom, RiseUpp get annually 2000 SEK and 20 SEK per member. RiseUpp aims to be an association at Uppsala university that gathers all students interested in space, hosts events such as lunch lectures, and participate in programmes like Rexus Bexus.

The \gls{mirage} project will consist mostly of students working on their thesis, therefore Uppsala university covers costs up to 1000 SEK per student. In addition, through the team's supervisor, Uwe Zimmermann, access to a laboratory with the necessary tools and a vacuum chamber for testing is given. Uppsala university also hosts the Swedish Institute for Space Physics (Institutet för Rymdfysik), which has a thermal chamber that will be used for testing.

RiseUpp has an agreement with the Swedish company SenseAir AB, in which SenseAir AB has provided a brand new K96 sensor, a used K96 sensor for testing, \gls{ir} lamps (which are a critical part of the K96 sensor) and will provide extra gas sensors if needed.

In addition, the sensor calibrations will be done in collaboration with the Norunda weather stations; who will provide gas samples and a Picarro gas analyzer for reference.

The gas pumps and the compressor are purchased and received, which is another crucial component to the mission.

% %%%%%%%%%%%%%%%%%%%%%%%%%%%%%%%%%%%%%%%%%%%%%%%%%%%%%%%%%%%%%%%%%%%%%%%%%%%%%%%
\subsection{Workforce}

The workforce consists of voluntary and thesis students. Voluntary students spend around 10 hours a week and will be the main workforce until March, they do this voluntarily and will be the core team members. Thesis students are recruited to begin working on the project from March, they will work in project groups under the supervision of the team's supervisor, Associate Professor Uwe Zimmermann, and will get credited for the work. The thesis project will be 15 credits, and the students are expected to work full-time on the project from 23rd of March until the 5th of June.

More information about the thesis projects is shown in \tref{tab:thesis}. Thesis students will not directly write into the \gls{sed} but they will write their own university reports.

One task for the voluntary students is to mentor and follow up with the thesis students. The voluntary students cover all subteams and will act as the interface between the thesis projects and the \gls{mirage} project.

There are weekly meetings for the whole team in which the core team members will follow up closely with the thesis projects and update on the progress, in addition to an extra subteam meeting each week to work on given tasks and plan the next steps. Thesis students will have the chance, and be encouraged, to continue working on the project voluntarily.

\begin{table}[htbp]
\centering
\caption{Thesis Projects Overview}
\label{tab:thesis}
\begin{tabular}{|l|L{3.5cm}|L{5.5cm}|}
\hline
\textbf{Project Number} & \textbf{Project Name} & \textbf{Project Description} \\ \hline

1 & Finite Element Analysis (FEA) and Computational Fluid Dynamics (CFD) 
& Finite Element Analysis (FEA) of temperature and stress \& Computational Fluid Dynamics (CFD) to analyze airflow. 
\\ \hline

2 & PCB Design and Embedded Programming 
& Printed circuit board (PCB) design, component selection, and embedded programming. 
\\ \hline

3 & Pressurized Gas Sensing Compartment 
& Design and assembly of a pressurized gas sensing experiment compartment. 
 \\ \hline

4 & Data Analysis Pipeline 
& Data analysis pipeline programming. 
 \\ \hline

5 & NDIR Sensor Testing and Calibration 
& Sensor testing and calibration of an NDIR gas sensor for stratospheric pressure and temperature conditions. 
 \\ \hline

6 & Thermal System Design 
& Thermal system design. 
 \\ \hline

7 & ECSS-Based Material Selection 
& Material selections in accordance with ECSS guidelines for a leakage-free and heat-resistant compartment. 
 \\ \hline

8 & Power Distribution Board Design 
& Power distribution board design. 
\\ \hline

\end{tabular}
\end{table}

In addition to the \gls{mirage} workforce, there is the RiseUpp board which will ensure the consistency of the team and assemble a team for the next Rexus Bexus cycles.
\clearpage

% %%%%%%%%%%%%%%%%%%%%%%%%%%%%%%%%%%%%%%%%%%%%%%%%%%%%%%%%%%%%%%%%%%%%%%%%%%%%%%%
\subsection{Budget}

The estimated total budget of the experiment \emph{parts} without the K96 sensor is 6089 SEK, and with a 15\% buffer is 7000 SEK ($\approx$ 650 EUR). The K96 sensor, which was sponsored, is not commercially available and was estimated to be 5000 SEK; making the total cost of the experiment around 1200 EUR. Additional managerial costs like outreach and travel expenses lead to a sum of approximately 85000 SEK ($\approx$ 8000 EUR).

For the hardware, \gls{mirage} was sponsored by SenseAir AB and received 2 K96 sensors; one for testing and one for use. Additionally, the pumps and the compressor were ordered.

All of the hardware costs can be covered by the money given through Uppsala university for the thesis projects, however \gls{mirage} aims to find external sponsors for additional outreach budget. Additional material such as cables, screws, RTCs and small microcontrollers are available for the team in the designated lab at Uppsala university and is provided; they were not added to the total cost.

\begin{longtable}{p{3cm}p{1.7cm}p{1.5cm}p{3cm}p{5cm}}
\caption{Budget Breakdown}
\label{tab:budget}\\

\toprule
\textbf{Cost Name} & \textbf{Quantity} &
\textbf{Unit Price (SEK)} &  \textbf{Status} & \textbf{Description} \\
\midrule
\endfirsthead

\multicolumn{5}{c}{\tablename\ \thetable\ -- \textit{Continued from previous page}}\\
\toprule
\textbf{Cost Name} & \textbf{Quantity} &
\textbf{Unit Price (SEK)} &  \textbf{Status} & \textbf{Description} \\
\midrule
\endhead

\midrule
\multicolumn{5}{r}{\textit{Continued on next page}}\\
\endfoot

\bottomrule
\endlastfoot

Enclosure & 1 & 720 & Not ordered &  \\
MS5607 Pressure Sensor & 2 & 81.75 & Not ordered & \\
SHT45 Temperature Sensor & 3 & 51.88 & Not ordered & \\
Voltage Regulators & 3 & 65.29 & Not ordered & \\
SD Card & 1 & 50.04 & Not ordered & 16 GB \\
\gls{pcb} & 1 & 200 & Not ordered & \\
Vacuum Pump & 2 & 1353.91 & Ordered & \\
Compressor & 1 & 1169.77 & Ordered & \\
Pump DevBoard & 1 & 53.96 & Ordered & \\
Power connector & 1 & 80.55 & Not ordered & \\
Amphenol Data Connector (RJF21B) & 1 & 571.25 & Not ordered & \\
10 W heating pad & 1 & 20.63 & Not ordered & \\
K96 & 2 & 5000 & Available & Non-commercial, estimate of value\\
\hline
15\% Margin for Hardware & & 913 & & \\
Total for Hardware & & 12000 & & \\
\hline
Travel Expenses & 5 & 3000 & Not ordered & For the launch campaign, counting only additional nonsponsored students \\
Accommodation & 5 & 10000 & Not ordered & For the launch campaign, counting only additional nonsponsored students \\
Team Building & 15 & 500 & Not ordered & \\
Team Jersey & 10 & 100 & Not ordered & \\
\hline
15\% Margin for Manegerial Costs & & 11025 & & \\
Total for Manegerial Costs & & 84500 & & \\
\end{longtable}


%
%\begin{table}[H]
%\centering
%\caption{Budget Breakdown}
%\label{tab:budget}
%\begin{tabularx}{\textwidth}{p{3cm}p{1.7cm}p{2cm}p{1cm}X}
%\toprule
%\textbf{Cost Name} & \textbf{Quantity} & \textbf{Category} & \textbf{Unit Price (SEK)} & \textbf{Description} \\
%\midrule
%Aluminum Block & 1 & Mechanical & 200 & For enclosure \\
%Plexy Glass & 1 & Mechanical & 90 & For lid \\
%Screws & 8 & Mechanical & 1 & \\
%High Precision Valve & 1 & Mechanical & 1000 & \\
%Cables & 20 & Electrical & 0.2 & \\
%MS5607 Pressure Sensor & 1 & Electrical & 81.75 & \\
%Temperature Sensor & 1 & Electrical & 93.74 & \\
%SHT3x-DIS Humidity Sensor & 1 & Electrical & 103.01 & \\
%Arduino Ethernet Shield & 1 & Electrical & 369 & \\
%Arduino Mega & 1 & Electrical & 579 & \\
%Voltage Regulators & 3 & Electrical & 60 & \\
%SD Card & 1 & Electrical & 119 & 32 GB \\
%\gls{pcb} & 1 & Electrical & 200 & \\
%Pump & 1 & Electrical & 500 & \\
%Thermal System & 1 & Electrical & 271.3 & \\
%Capacitor & 3 & Electrical & 1 & \\
%K96 & 1 & Electrical & 5000 & Non-commercial, estimate of value \\
%Travel Expenses & 5 & Managerial & 3000 & For the launch campaign, counting only additional nonsponsored students \\
%Accomodation & 5 & Managerial & 10000 & For the launch campaign, counting only additional nonsponsored students \\
%Team Building & 15 & Managerial & 500 & \\
%Team Jersey & 10 & Managerial & 100 & \\
%\bottomrule
%\end{tabularx}
%\end{table}

% %%%%%%%%%%%%%%%%%%%%%%%%%%%%%%%%%%%%%%%%%%%%%%%%%%%%%%%%%%%%%%%%%%%%%%%%%%%%%%%
\subsection{External Support}

\gls{mirage} is endorsed and adviced by the Uppsala university department of meteorology. In addition, \gls{mirage} is in the process of becoming an associated association to the student union for science and engineering at Uppsala university (UTN). This will open up the doors to apply for research grants provided by the UTN. Astronomisk Ungdom has a research grant pool of 100000 SEK which is available for associations under their umbrella, which will be applied to.

\gls{mirage} is in talks with Uppsala University Innovation and the development office at the university about financial and outreach support. They will connect \gls{mirage} with companies that work in the field and help with social events organized by RiseUpp.

% %%%%%%%%%%%%%%%%%%%%%%%%%%%%%%%%%%%%%%%%%%%%%%%%%%%%%%%%%%%%%%%%%%%%%%%%%%%%%%%

\section{Outreach Approach}
\subsection{Goal of the outreach}
The aim is to spread the word of the Rexus/Bexus programme and MIRAGE to as large an audiance as possible. Especially inspiring STEM majors at Uppsala university to engage in MIRAGE or similar Rexus/Bexus projects and experiments in the future.
\vspace{3mm}\\
The MIRAGE team has advertised themselves by posting continuous updates on social media about the experiment, team meetings, team building activities and humorous content. They have also written an informative letter about the Rexus Bexus programme, the RiseUpp association and MIRAGE that has been announced on several engineering program's studium pages for recruitment of bachelor thesis students.
\vspace{3mm}\\
Other media sources that the team has advertised themselves via, are newspapers. An article was written in the students university´s magazine Techna.
\vspace{3mm}\\
To cover equipment costs for the experiment and travel expenses the team has contacted local and national STEM-companies, authorities and foundations to ask for sponsorships and collaborations. One foundation that has been contacted was Sparbanksstiftelsen Upland, Swedbank and the team have an ongoing collaboration with RedBull.

\begin{table}[H]
\centering
\caption{Team social media links}
\label{tab:social_media}
\begin{tabular}{|c|c|c|c|}\hline
\textbf{Team} & \textbf{Instagram} & \textbf{LinkedIn} & \textbf{Website} \\\hline
Mirage & \href{https://www.instagram.com/riseupp_uu/?igsh=MXhsM3lsZmMyMjh4ZA\%3D\%3D#}{\text{riseupp\_uu}} & \href{https://www.linkedin.com/company/riseupp-uu/posts/?feedView=all}{riseupp-uu}&  https://www.spacestudents.se/riseupp/   \\\hline
\end{tabular}
\end{table}

\subsection{Future outreach plans}
In the future, the team plans to arrange lunch lectures with SSC, Christer Fuglesang, Marcus Wandt. These events can be sponsored by Uppsala kommun and supported by RedBull to attract a greater audience.
\vspace{3mm}\\
The team will continue to stay in touch with companies, authorities and foundations as well as contacting new potential sponsors.
\vspace{3mm}\\
The team wants to reach out to other media sources such as Uppsala Nya Tidning, Sveriges Radio and Uppsala university's student website/social media.
\vspace{3mm}\\
In a later stage of the project, a documentary movie will be released that follows the team´s project journey and the launch from Kiruna in October 2026.



\section{Risk Register}

\textbf{Risk ID}

TC -- technical/implementation

MS -- mission (operational performance)

SF -- safety

VE -- vehicle

PE -- personnel

EN -- environmental

OR -- organisational

DS -- data related

\textbf{Probability (P)}

A. Minimum -- Almost impossible to occur

B. Low -- Small chance to occur

C. Medium -- Reasonable chance to occur

D. High -- Quite likely to occur

E. Maximum -- Certain to occur, maybe more than once

\textbf{Severity (S)}

1. Negligible -- Minimal or no impact

2. Significant -- Leads to reduced experiment performance

3. Major -- Leads to failure of subsystem or loss of flight data

4. Critical -- Leads to experiment failure or creates minor health hazards

5. Catastrophic -- Leads to termination of the REXUS and/or BEXUS programme, damage to the vehicle or injury to personnel

The rankings for probability (P) and severity (S) are combined to assess the overall risk classification, ranging from very low to very high and being coloured green, yellow, orange or red according to the \gls{sed} guidelines.

\begin{longtable}{L{1.5cm}L{3cm}L{1cm}L{1cm}L{2cm}L{6cm}}
\caption{Risk Register}
\label{tab:risk_register}\\

\toprule
\textbf{ID} & \textbf{Risk \& consequence} & \textbf{P} & \textbf{S} &
\textbf{P x S} & \textbf{Prevention and/or Mitigation Action} \\
\midrule
\endfirsthead

\multicolumn{6}{c}{\tablename\ \thetable\ -- \textit{Continued from previous page}}\\
\toprule
\textbf{ID} & \textbf{Risk \& consequence} & \textbf{P} & \textbf{S} &
\textbf{P x S} & \textbf{Prevention and/or Mitigation Action} \\
\midrule
\endhead

\midrule
\multicolumn{6}{r}{\textit{Continued on next page}}\\
\endfoot

\bottomrule
\endlastfoot

\rowcolor{mediumorange} TC10 & Pump failure & C & 4 & medium & Have multiple pumps in setup, choose appropriate components \\
\rowcolor{lowyellow} TC11 & Pump failure & C & 2 & low &  Acceptable risk \\
\rowcolor{mediumorange} MS10 & Over-frozen gas inlet or outlet & C & 4 & medium & Large enough radius of outlets and inlets, heating element at inlet to heat air to appropriate temperature, move primary expansion point of air to being inside the experiment enclosure \\
\rowcolor{verylowgreen} MS11 & Over-frozen gas inlet or outlet & A & 4 & very low & Acceptable risk \\
\rowcolor{lowyellow} DS10 & Loss of data & B & 4 & low & Send part of data to ground regularly \\
\rowcolor{verylowgreen} DS11 & Loss of data & A & 4 & very low & Acceptable risk \\
\rowcolor{mediumorange} TC20 & Underheating of electronics during ascent & C & 3 & medium & More heating measures such as insulation and heating elements \\
\rowcolor{verylowgreen} TC21 & Underheating of electronics during ascent & A & 4 & very low & Acceptable risk \\
\rowcolor{lowyellow} TC30 & Failure of thermal control due to too low temperatures & B & 3 & low & Acceptable risk\\
\rowcolor{verylowgreen} PE10 & Drop out of Team member & B & 2 & very low & Acceptable risk \\
\rowcolor{mediumorange} OR10 & Lack of funding & C & 4 & medium & Contact more companies, find other sources from university \\
\rowcolor{lowyellow} OR11 & Lack of funding & B & 4 & low & Acceptable risk  \\
\rowcolor{lowyellow} MS20 & \gls{mcu} failure & B & 4 & low & Test under pressure and temperature conditions while on ground and choose different component if it doesn't tolerate, use watchdog timer and consider redundant systems\\
\rowcolor{verylowgreen} MS21 & \gls{mcu} failure & A & 4 & very low & Acceptable risk\\
\rowcolor{verylowgreen} TC40 & Outside sensors fail & B & 2 & very low & Acceptable risk \\
\rowcolor{lowyellow} TC50 & Condensation of water inside the measurement chamber & C & 2 & low & Implement flushing cycle \\
\rowcolor{verylowgreen} TC51 & Condensation of water inside the measurement chamber & B & 2 & very low & Acceptable risk \\
\rowcolor{lowyellow} TC60 & Contamination due to materials of the balloon, gondola, and experiment setup & C & 2 & low & Carefull selection of non-outgassing materials of the experiment and check with organisers for balloon/gondola materials; direct the inlet away from the ballon and other experiments\\
\rowcolor{verylowgreen} TC61 & Contamination due to materials of the balloon, gondola, and experiment setup & A & 2 & very low & Acceptable risk\\
\rowcolor{lowyellow} MS30 & Software failure of thermal control & B & 3 & low & Implement a watchdog and an option to restart software from ground. Include thermal switch to disengage heating devices if temperatures are too hot. \\
\rowcolor{verylowgreen} MS31 & Software failure of thermal control & A & 3 & very low & Acceptable risk \\
\rowcolor{lowyellow} MS40 & Software failure of pressure control & C & 2 & low & Implement watchdog and option to restart software from ground \\
\rowcolor{verylowgreen} MS41 & Software failure of pressure control & A & 2 & very low & Acceptable risk \\
\end{longtable}
% %
% %\begin{landscape}
% %\begin{table}[H]
% %\centering
% %\caption{Risk Register}
% %\label{tab:risk_register}
% %\begin{tabular}{p{1.5cm}p{4cm}p{1cm}p{1cm}p{2cm}p{6cm}}
% %\toprule
% %\textbf{ID} & \textbf{Risk \& consequence} & \textbf{P} & \textbf{S} & \textbf{P x S} & \textbf{Prevention and/or Mitigation Action} \\
% %\midrule
% %TC10 & Pump failure & C & 4 & medium & have multiple pumps in setup, chose appropriate components \\
% %MS10 & Over frozen gas inlet or outlet & C & 4 & medium & Large enough radius, heat air to appropriate temperature \\
% %DS10 & Loss of data & B & 4 & low & Send part of data to ground regularly \\
% %TC20 & Over/Underheating of electronics & C & 4 & medium & More heating/cooling measures depending on test results \\
% %TC30 & Failure of thermal control & C & 3 & low & Planning and testing \\
% %PE10 & Drop out of Team member & C & 2 & low & Good communication between team members about workload and availability; additional recruiting to reduce workload; selection for team members with sufficient time resources during application \\
% %OR10 & Lack of funding & B & 5 & medium & Contact companies soon \\
% %MS20 & \gls{mcu} failure & B & 4 & low & Test under pressure and temperature conditions while on ground \\
% %MS30 & Automatic switching of measurement times fails & C & 2 & low & Add manual option as backup \\
% %MS40 & Experiment doesn't respond to manual commands from ground & C & 1 & Very low & No action \\
% %MS50 & Loss of communication with ground stations & C & 1 & Very low & No action \\
% %TC40 & Outside sensors not working due to low temperatures & C & 2 & low & Chose and test appropriate components \\
% %TC50 & Condensation of water inside the \gls{ndir} sensor box. & C & 4 & medium & Add water removal system to design \\
% %DS20 & Data storage overflow & A & 4 & Very low & SD card should be chosen large enough with enough margin \\
% %MS60 & Software failure for thermal control & B & 3 & low & Option to restart software from ground \\
% %MS70 & Software failure for pressure control & B & 4 & low & Option to restart software from ground \\
% %\bottomrule
% %\end{tabular}
% %\end{table}
% %\end{landscape}
% %

% %%%%%%%%%%%%%%%%%%%%%%%%%%%%%%%%%%%%%%%%%%%%%%%%%%%%%%%%%%%%%%%%%%%%%%%%%%%%%%%\newpage

\chapter{EXPERIMENT DESCRIPTION}
\label{sec:experiment_description}
\section{Experiment setup}

\begin{figure}[htbp]
    \centering
    \includegraphics[width=0.8\textwidth]{0_images/figures/Mirage Block schematic.png}
    \caption{Experimental Setup}
    \label{fig:experimental_setup}
    
\end{figure}

The experiment will feature a NDIR sensor (K96), a pressurization system, a thermal system, ambient environment sensors and a processing unit. 
The pressurization system will include a pumping mechanism in order to increase the pressure and the amount of particles that can be measured. The air will be heated up and pumped into the measurement chamber with the K96 sensor inside together with temperature, pressure and humidity sensors. The gas will be exchanged continously to exchange the sample while maintaining the pressure of 3 bar. 
The K96 Sensor will communicate with the processing unit and the processing unit sends the measured digital data for storage to the SDHC card.
The pressure sensor,temperature sensor and humidity sensor communicate with the processing unit via digital interfaces (such as $I^2C$) to ensure that the pump and thermal control system are running with the right parameters to maintain room temperature conditions and a stable pressure of 3 bar. Unique addresses are used to identify each sensor's signal, allowing the processing unit to collect and organize data from all ambient sensors efficiently.

%%%%%%%%%%%%%%%%%%%%%%%%%%%%%%%%%%%%%%%%%%%%%%%%%%%%%%%%%%%%%%%%%%%%%%%%%%%%%%%
\section{Experiment Interfaces}

%%%%%%%%%%%%%%%%%%%%%%%%%%%%%%%%%%%%%%%%%%%%%%%%%%%%%%%%%%%%%%%%%%%%%%%%%%%%%%%
\subsection{Mechanical Interface}\label{mechanical_interface}

The \gls{mirage} enclosure is mechanically attached to the ESCARGO gondola structure via four mounting points located on its bottom plate. These four connection points interface with two parallel 45 $\times$ 45 mm Bosch aluminium T-slot profiles using M6 T-bolts, M6 rubber spacers, and M6 flange nuts. The parallel profiles are then mounted to the ESCARGO gondola structure using standard Rexroth M6 brackets, M6 T-bolts, and M6 flange nuts. This configuration facilitates installation without necessitating structural modifications to the ESCARGO frame, ensuring full compatibility with the standard mounting interface.

To reduce vibration transmission, temperature transfer and accommodate minor alignment tolerances during installation, rubber spacers are installed between the enclosure and the mounting brackets.

The mounting configuration provides a rigid mechanical interface between the experiment enclosure and the ESCARGO frame, ensuring structural redundancy and compliance with the standard ESCARGO mechanical interface requirements.

\begin{table}[htbp]
\centering
\begin{tabular}{lcl}
\toprule
\textbf{Item} & \textbf{Quantity} & \textbf{Specification} \\
\midrule
L-bracket & 4 & Bosch Rexroth aluminium M6 bracket \\
T-Bolt & 12 & M6 $\times$ 25 mm \\
Flange nut & 12 & M6 \\
Rubber Spacer/Damper & 4 & M6 \\
\bottomrule
\end{tabular}
\caption{Mounting hardware components}
\label{tab:mounting_hardware}
\end{table}

\begin{figure}[H]
    \centering
    \includegraphics[width=0.45\textwidth]{0_images/figures/mouting.png}
    \includegraphics[width=0.45\textwidth]{0_images/figures/render3.png}
    \caption{CAD images of the mechanical interface.}
    \label{fig:example_image}
    
\end{figure}
%%%%%%%%%%%%%%%%%%%%%%%%%%%%%%%%%%%%%%%%%%%%%%%%%%%%%%%%%%%%%%%%%%%%%%%%%%%%%%%
\subsection{Electrical Interface}\label{electrical_interface}

The electrical interface to the Bexus power unit consists of the MIL-C 26482P series 1 connector, as provided by the user manual. Using pin A and B, the supply voltage from the SAFT MP176065 battery pack will be fed to the power distribution board.
The connector to the E-Link of the gondola will be an Amphenol RJF21B and the RJ45 connector will be employed for communication with a RJ45 plug on the MCU PCB.

Since the Bexus power supply is 28.8V two LM2596HVS step-down voltage regulators will be used on the power distribution board. For the MCU and K96 sensor 8V will be used, the pressurization system operates at 24V, the thermal system at 12V and 3.3V for the sensors. The estimated average current draw from the BEXUS battery is 0.8A and the average power consumption is 23.11W. The total power consumption from BEXUS power unit is 138.67Wh as listed in section 4.7 and the peak current draw is 2.26A. For protecting the BEXUS batteries and other experiments a pre-charge circuit will be implemented in order to limit the currents of the pressurization system and an LC filter to deal with the EMI and switching frequencies of DC/DC converters. Distributed single point ground will be implemented.


%%%%%%%%%%%%%%%%%%%%%%%%%%%%%%%%%%%%%%%%%%%%%%%%%%%%%%%%%%%%%%%%%%%%%%%%%%%%%%%
\subsection{Thermal}\label{thermal}
As outlined in \sref{sec: thermal design}, the thermal behaviour of the experiment will be controlled via heating elements and insulation. To avoid thermal coupling to ESCARGO, rubber bumpers between the experiment and the gondola will be used.
The thermal coupling will be further tested when the experiment is assembled. The air inlet and outlet for the air samples to be measured will be directed away from the vehicle and should not heat or cool the vehicle.

%%%%%%%%%%%%%%%%%%%%%%%%%%%%%%%%%%%%%%%%%%%%%%%%%%%%%%%%%%%%%%%%%%%%%%%%%%%%%%%
\section{Experiment Components}

\begin{table}[H]
\centering
\caption{Experiment summary table}
\label{tab:experiment_summary}
\begin{tabular}{L{6cm}L{6cm}}
\toprule
Experiment mass (in kg): & $\sim$4.0 \\
Experiment dimensions (in m): & 0.20x0.20x0.15\\
Experiment expected \gls{cog} position: & Distance from geometric centre: x=4mm y=23mm z=22mm \\
\bottomrule
\end{tabular}
\end{table}

\begin{table}[H]
\centering
\caption{Electrical components}
\label{tab:electrical_components}
\begin{tabular}{p{4cm}p{4cm}ll}
\hline
\textbf{Component} & \textbf{Dimensions (mm)} & \textbf{Mass (g)} & \textbf{Quantity}\\
\hline
Vacuum pump & 119$\times$44$\times$73 & 294 & 2 \\
Compressor & 130$\times$60$\times$85 & 445 & 1 \\
Arduino Mega & 54$\times$102$\times$15 & 37 & 1\\
K96 & 42.26$\times$33.9$\times$26.04 & 300 & 1 \\
MS5607 & 3.0$\times$5.0$\times$0.95 & <1 & 1\\
SHT45 & 1.5$\times$1.5$\times$0.5 & <1 & 4 \\
LM2596HVS-ADJ & 10.1$\times$9.0$\times$4.5 & 1.59 & 3 \\
\gls{pcb} & TBD & TBD & 1 \\
Thermal system & TBD & TBD &  \\
\hline
Total & & 1079.59 & \\
\hline
\end{tabular}
\end{table}

%%%%%%%%%%%%%%%%%%%%%%%%%%%%%%%%%%%%%%%%%%%%%%%%%%%%%%%%%%%%%%%%%%%%%%%%%%%%%%%
\section{Mechanical Design}\label{mechanical_design}

The mechanical design of the \gls{mirage} experiment focuses on a compact enclosure housing a pressurized measurement chamber, a multi-stage pump system, supporting electronics, and a heating system. \gls{mirage} is designed to fit within the ESCARGO gondola structure while maintaining structural integrity and ensuring safe operation under the environmental conditions expected during the mission.

\begin{figure}[htbp]
    \centering
    \includegraphics[width=0.8\textwidth]{0_images/figures/tech_drawing.png}
    \caption{Technical drawing of the system as a whole.}
    \label{fig:tech_drawing}
    
\end{figure}

\subsection{Experiment Enclosure}

The experiment is housed in a \gls{cots} aluminium enclosure with external dimensions of $200\times 200\times 150$ mm and a wall thickness of expected 2 mm. The enclosure provides mechanical protection for all internal components and acts as the primary structural body of the experiment. A \gls{cots} enclosure was selected to reduce manufacturing complexity, minimize cost, and provide a robust structural housing with well-characterized mechanical properties.
The enclosure is directly mounted to the ESCARGO frame as described in \sref{mechanical_interface}. Internally, the enclosure accommodates the pressure chamber, two diaphragm vacuum pumps, one piston compressor, a passive outflow hole, shutter valve, tubing, and electronic control hardware. The enclosure also includes two dedicated holes for air inlet and one for outlet, as well as electrical feedthroughs required for connecting to the BEXUS power and E-link system.

\begin{figure}[htbp]
    \centering
    \includegraphics[width=0.45\textwidth]{0_images/figures/tech_drawing_3.png}
    \includegraphics[width=0.45\textwidth]{0_images/figures/tech_drawing_5.png}
    \caption{Technical drawings of experiment enclosure. Drawings are representative and the thickness of both “Main Body” and “Main Lid” will be reduced to 2mm by buying COTS parts.}
    \label{fig:enclosure}
\end{figure}

\subsection{Internal Layout}

The internal layout of the \gls{mirage} experiment is arranged to maintain a compact configuration and minimize the length of the airflow path between the pressurized chamber and the compressor.
Electronic control hardware is mounted along the inner bottom of the enclosure, providing physical isolation between electrical components and the pressurized air system.
This configuration ensures efficient airflow through the chamber while maintaining accessibility for integration and servicing.

\begin{figure}[htbp]
    \centering
    \includegraphics[width=0.8\textwidth]{0_images/figures/tech_drawing_6.png}
    \caption{Drawing of the internal layout.}
    \label{fig:int_layout}
    
\end{figure}

\subsection{Pressurized Measurement Chamber}

The core element of the experiment is a small pressurized chamber that houses the K96 sensor used for the measurements. The chamber has internal dimensions of approximately $38 \times 33 \times 45$ mm and a wall thickness of 7 mm. It will be machined from aluminium by the \gls{mirage} team at Uppsala University.

The chamber is sealed using a copper gasket to ensure reliable sealing under elevated pressure conditions and temperature changes. Copper gaskets were selected due to their ability to provide a robust metal-to-metal seal capable of maintaining integrity at high pressure differentials.

The chamber is designed to minimize internal volume while providing sufficient space for the sensor and the necessary airflow path. Reducing the internal volume proportionally lowers the airflow rate required to exchange the chamber atmosphere, which improves the overall performance of the experiment.

\begin{figure}[htbp]
    \centering
    \includegraphics[width=0.45\textwidth]{0_images/figures/render2.png}
    \includegraphics[width=0.45\textwidth]{0_images/figures/render1.png}
    \caption{Renders of the pressurized chamber.}
    \label{fig:pressure_render}
\end{figure}

\begin{figure}[htbp]
    \centering
    \includegraphics[width=0.45\textwidth]{0_images/figures/tech_drawing_2.png}
    \includegraphics[width=0.45\textwidth]{0_images/figures/tech_drawing_4.png}
    \caption{Technical drawings of the pressurized chamber.}
    \label{fig:pressure_drawing}
\end{figure}

\subsection{Pump System Architecture}\label{sec: pump system arch}

The \gls{mirage} experiment utilizes a multi-stage air circulation system to maintain a controlled internal pressure while ensuring a continuous exchange of air within the measurement chamber.

\textbf{System Operation:}

\begin{enumerate}
    \item Pre-heating: Ambient intake air is pre-heated before entering the vacuum stage to prevent condensation or freezing within the system components.
    \item Vacuum Stage: Two diaphragm vacuum pumps operate continuously from launch up to an altitude of 22 km. These pumps supply the secondary stage with air at approximately atmospheric pressure ($\approx$ 1 bar). Since the volumetric flow rate decreases with decreasing intake pressure, the pumps will operate at higher settings at higher altitudes.

    \item Intercooling: To manage thermal loads generated by compression, the tubing connecting the vacuum pumps to the piston compressor serves as a passive intercooler.
    \item Compression Stage: A piston compressor increases the airflow pressure to a maximum of 3 bar before it enters the measurement chamber.
    \item Pressure Regulation: The internal chamber pressure is maintained by a passive outflow valve equipped with a shutter mechanism. This valve regulates the exhaust flow rate to the external environment.
    \item Monitoring: Integrated pressure and temperature sensors are situated within the tubing between compression stages to provide real-time performance data and ensure stable operation.
\end{enumerate}

%% ADD FIGURE %%

\textbf{Compressor Design Considerations:} Two primary design methodologies are being evaluated for the compression stage:

\begin{enumerate}
    \item Continuous Low-Flow Operation: Utilizing a compressor capable of maintaining a continuous, very low flow rate to match the chamber's exchange requirements.
    \item Intermittent Buffered Operation: Utilizing a higher flow rate compressor operated in duty cycles. This configuration requires a 1 bar buffer tank between pressurisation stages to accumulate air from the vacuum stage. The compressor activates to deplete the tank and deactivates once the tank pressure reaches a lower threshold.
\end{enumerate}

The nominal operating pressure inside the chamber is expected to reach approximately 3 bar above perfect vacuum. At float altitude, the external ambient pressure may decrease to approximately 10 mbar, resulting in a pressure differential of approximately 3 bar across the chamber walls. This pressure differential is accounted for in the structural design of the chamber.

The pressure between the vacuum pumps and the compressor is expected to remain close to atmospheric pressure ($\approx$ 1 bar) during normal operation in order to match expected compressor efficiency.

Temperature conditions at float altitude will require the implementation of pre-heating of the intake air before the vacuum stage to prevent condensation or freezing within the system. Additionally, pressure sensors will be placed between compression stages to monitor system performance and ensure stable pressure control during operation.

% \subsection{In-line Sensing and Environmental Monitoring}

% To ensure the operational integrity of the gas handling system and to facilitate real-time performance analysis, the \gls{mirage} experiment incorporates an integrated sensing suite within the tube circuit. This monitoring architecture is essential for validating the state of the air as it transitions through the various compression and thermal management stages.

% \begin{enumerate}
%     \item Pressure Sensing: High-precision pressure sensors are situated within the tubing between the diaphragm vacuum pumps and the piston compressor. These sensors monitor the inter-stage pressure to ensure that the vacuum pumps, operating continuously up to 22 km, maintain the required 1 bar supply pressure for the compressor.
%     \item Temperature Sensing: Thermal sensors are embedded directly into the flow path to monitor the effectiveness of the pre-heating and intercooling mechanisms. Monitoring the air temperature post-heating is critical to verify that the intake air has reached a safe threshold before entering the pump heads and K96 sensor chamber, thereby preventing damage from cold or heat
% \end{enumerate}

\subsection{Pump and Fluid Handling Technologies}

Small gas circulation and sampling systems used in atmospheric sensing instruments typically rely on compact electrically driven pumps to move air through measurement chambers and sensors. Several pump technologies are commonly used in portable gas analysis systems, including diaphragm pumps, piston compressors, rotary vane pumps, and scroll pumps. Each technology offers different advantages in terms of pressure capability, size, reliability, and maintenance requirements.

Diaphragm pumps are widely used in environmental monitoring and portable gas sampling instruments. These pumps operate by oscillating a flexible diaphragm using an electric motor, creating alternating suction and compression cycles. Because the gas is hermetically separated the mechanical drive components by the diaphragm, diaphragm pumps operate oil-free and are resistant to contamination. This makes them particularly suitable for experiments involving sensitive sensors or atmospheric gas sampling. Furthermore, diaphragm pumps are compact, lightweight, and tolerant of condensation in the airflow, which is expected in the \gls{mirage} experiment due to low external temperatures at float altitude.

Piston compressors are commonly used when higher pressures are required. These pumps use a reciprocating piston to compress air within a cylinder, allowing significantly higher pressure ratios than typical diaphragm pumps. However, piston compressors generally contain more moving mechanical parts and may be more sensitive to contamination or thermal effects. For this reason, piston compressors are often used in series with diaphragm pumps in small air circulation systems, where the diaphragm pumps provide initial airflow and the compressor increases the pressure to the required level.

Other pump technologies, such as rotary vane pumps and scroll pumps, are frequently used in laboratory vacuum systems. While these technologies can achieve significantly lower pressures and higher pumping speeds, they typically require lubrication, larger mechanical assemblies, and higher power consumption. Due to these characteristics, they are generally unsuitable for compact balloon payload experiments such as \gls{mirage}, where mass, power consumption, and system simplicity are critical design constraints.

Based on this comparison, diaphragm pumps were selected as the primary airflow drivers in the \gls{mirage} experiment due to their oil-free operation, compact size, reliability, and resistance to condensation. A piston compressor is used downstream of the diaphragm pumps to achieve the required pressurization of the measurement chamber.

\subsection{Tubing and Flow Control Components}

Gas flow regulation in compact atmospheric sensing instruments is traditionally achieved through the use of miniature solenoid or needle valves. Solenoid valves are often preferred for their rapid response times and ability to provide high-precision, electronically controlled flow regulation within automated systems.
However, for the \gls{mirage} experiment, the design has been optimized to utilize a passive outflow valve equipped with a mechanical shutter. This shift to a passive regulation system is engineered to maintain the required internal chamber pressure of approximately 3 bar by modulating the exhaust flow rate (hole size) based on the pressure differential encountered during the flight profile. In addition the usage of a shutter-based passive mechanism rather than an active solenoid, the system achieves greater mechanical simplicity and safety if the ballon would fall into the water.
The final selection of the hole and the shutter's mechanical properties is based on ensuring a stable airflow path that accommodates the continuous operation of the upstream diaphragm pumps up to the 22 km altitude threshold.

\subsection{Tubing and Fluid Connections}

The airflow between the pumps, compressor, and pressurized measurement chamber is transported through flexible tubing. In compact gas circulation systems, tubing connections are typically achieved using barbed fittings, compression fittings, or push-to-connect pneumatic fittings.

Barbed fittings are commonly used with flexible tubing materials such as silicone or polyurethane, where the tubing is pressed onto the fitting and secured through friction or hose clamps. Compression fittings provide a more robust and leak-tight connection and are frequently used in higher-pressure gas systems. Push-to-connect fittings are widely used in pneumatic systems due to their ease of installation and reliable sealing.

For the \gls{mirage} experiment, the tubing connections will be selected to ensure leak-tight operation at pressures up to approximately 3 bar while maintaining compatibility with the chosen tubing material and the compact geometry of the experimental setup. Particular attention will be given to minimizing leakage and dead volumes in the system to maintain stable pressure control and reliable airflow through the measurement chamber.

In alignment with the experiment's thermal management requirements, the intake tube lines will be integrated with high-conductivity copper windings to facilitate efficient heat transfer between the enclosure and the vacuum pump interface. This thermal exchange mechanism is designed to rapidly elevate the temperature of the intake air, ensuring the fluid stream remains within the pump's safe operational range. This will also be the case between compressor and the chamber as the heat created by compression could affect the quality of the study. This proactively mitigates the risk of mechanical failure or structural damage to the diaphragm components caused by the extreme low-temperature environment encountered during the ascent to 22 km and high temperatures after compression.

\subsection{Pumps}

The air circulation system utilizes commercially available miniature medical pumps. The vacuum pumps are compact DC diaphragm pumps capable of generating the suction required to draw air into the system at 10 mbar. Two pumps are used to provide sufficient airflow at 1 bar for the compressor and to enhance system robustness.

The compressor is a small DC piston air pump used to increase the pressure of the air entering the pressurized chamber. Both components were selected based on their cost, compact size, low mass, and compatibility with the power constraints of the BEXUS payload system.

\subsection{Structural Considerations}

The mechanical structure is designed with sufficient strength and safety margin to withstand internal pressure differentials exceeding the maximum expected value of approximately 3 bar between the pressurized chamber and the external environment.
The pressurized chamber is manufactured from solid machined aluminium with a wall thickness of 7 mm. This thickness provides a structural safety factor far greater than 2 relative to the maximum expected pressure differential, accounting for uncertainties in material properties, manufacturing tolerances, and thermal stresses.
The outer aluminium enclosure provides mechanical protection and structural support for the internal components. A copper gasket is used to ensure reliable sealing and minimize the risk of leakage during operation.
The internal layout maintains a compact configuration while ensuring that mechanical loads generated by the pumps, tubing, and mounting hardware are safely transferred to the enclosure structure.

\subsection{Mass Budget}

A preliminary mass budget of the \gls{mirage} experiment was derived from the SolidWorks CAD assembly model. The estimated mass distribution of the main structural, fluid handling, electronic, and fastening components is summarized in \tref{tab:mass_budget}. Smaller components such as fasteners and connectors are included to provide a realistic estimate of the total system mass.

\begin{table}[htbp]
\centering
\begin{tabular}{lccc}
\toprule
\textbf{Component} & \textbf{Qty} & \textbf{Mass (g)} & \textbf{Total (g)} \\
\midrule

\multicolumn{4}{l}{\textbf{Structural components}} \\

COTS aluminium enclosure & 1 & 1200 (TBD) & 1200 \\
Pressurized chamber body & 1 & 159 & 159 \\
Pressurized chamber lid & 1 & 70 & 70 \\
Vacuum pump arch & 1 & $\sim$30 & 30 \\
Valve bracket & 1 & $\sim$20 & 20 \\

\addlinespace
\multicolumn{4}{l}{\textbf{Fluid system}} \\

Diaphragm vacuum pump & 2 & $\sim$294 & 588 \\
Piston compressor & 1 & $\sim$445 & 445 \\
Shutter valve & 1 & 250 (TBD) & 250 \\
Air tubing assembly & 1 & $\sim$120 & 120 \\

\addlinespace
\multicolumn{4}{l}{\textbf{Electronics}} \\

Arduino Mega 2560 & 1 & 37 & 37 \\
RJF21B connector & 1 & $\sim$15 & 15 \\

\addlinespace
\multicolumn{4}{l}{\textbf{Fasteners}} \\

M3 screws & 13 & $\sim$1 & 13 \\
M3.5 screws & 4 & $\sim$2 & 8 \\
M3 long screws & 4 & $\sim$2 & 8 \\
M6 mounting bolts & 4 & $\sim$10 & 40 \\

\midrule
\textbf{Estimated total} & & & \textbf{$\sim$3003} \\
\bottomrule
\end{tabular}
\caption{Mass budget of system components.}
\label{tab:mass_budget}
\end{table}

\subsection{Component Procurement Summary}

The experiment design utilizes a combination of \gls{cots} components and custom-manufactured parts. \gls{cots} components are standard products purchased directly from manufacturers without modification, while custom components are specifically designed and machined for the \gls{mirage} experiment.

The \gls{cots} components include the aluminium enclosure, diaphragm vacuum pumps, piston compressor, shutter valve, tubing, and standard fasteners. These components were selected in order to reduce manufacturing complexity, cost, and lead time.

Custom-manufactured components include the pressurized measurement chamber (chamber body and lid) and several structural mounting elements such as pump holders, valve holders, and supporting brackets. These parts are machined from aluminium at Uppsala University to meet the structural and pressure requirements of the experiment.

The classification of the main components as \gls{cots} or custom-manufactured is summarized in \tref{tab:component_type_summary}.

\begin{table}[htbp]
\centering
\begin{tabular}{ll}
\toprule
\textbf{Component} & \textbf{Type} \\
\midrule

Main enclosure & COTS \\
Diaphragm vacuum pumps & COTS \\
Piston compressor & COTS \\
Shutter valve & COTS \\
Arduino Mega 2560 & COTS \\
RJF21B connector & COTS \\

Pressurized chamber body & Custom machined \\
Pressurized chamber lid & Custom machined \\
Valve holder & Custom machined \\
Pump mounting brackets & Custom machined \\
Structural holders and supports & Custom machined \\

Tubing and pneumatic fittings & COTS \\
Fasteners (M3/M6 bolts, screws) & COTS \\

\bottomrule
\end{tabular}
\caption{Component Type Summary}
\label{tab:component_type_summary}
\end{table}

%%%%%%%%%%%%%%%%%%%%%%%%%%%%%%%%%%%%%%%%%%%%%%%%%%%%%%%%%%%%%%%%%%%%%%%%%%%%%%%
\section{Electronics Design}

The overall electronics architecture of the \gls{mirage} experiment is shown in \fref{fig:electrical_schematic}. The system consists of a Power Distribution Board, which interfaces with the BEXUS gondola power supply and provides regulated voltages to the experiment subsystems. The processing core is based on an Arduino Mega microcontroller, which manages sensor data acquisition, system control, and communication with onboard storage.

Additional subsystems include the pressurization system, the thermal control system, the K96 NDIR sensor core, and external environmental sensors measuring ambient temperature, pressure, and humidity. The ambient sensors communicate with the processing unit via an I²C interface, while the K96 NDIR sensor communicates using a UART interface. The interactions between these subsystems and the processing unit are illustrated in \fref{fig:electrical_schematic}. Each subsystem is described in more detail in the following sections.

\begin{figure}[H]
    \centering
        \includegraphics[width=1.0\textwidth]{0_images/figures/Mirage_Electrical_Schematic.png}
            \caption{Electrical Schematic}
                \label{fig:electrical_schematic}
                \end{figure}

\subsection{Processing core unit}

%\begin{figure}[H]
%    \centering
%    \includegraphics[width=0.5\textwidth]{0_images/figures/Arduino MEGA Rev3.png}
%    \caption{Arduino Mega Rev3 WiFi}
%    \label{fig:component_arduino}
%\end{figure}


An Arduino Mega will be used for preliminary testing and as a backup solution, and will be replaced by a dedicated processing \gls{pcb} once assembly begins. The Arduino Mega is equipped with an Ethernet shield, enabling it to communicate with EGSE and save measurement data to an SD card. The Arduino needs a power input between 7-12V, which will be converted from the 28.8V provided by Bexus' SAFT batteries through a DC/DC stepdown regulator.

\subsection{MicroSD Card}

%\begin{figure}[H]
%    \centering
%    
\includegraphics[width=0.3\textwidth]{0_images/figures/MicroSD card.png}
%    \caption{SDHC Memory Card}
%    \label{fig:component_SDHC}
%\end{figure}

Measurement data will be stored on a microSDHC memory card (SP016GBSTHBU1V10SP) with a capacity of 16~GB. The card has maximum read and write speeds of up to 85~MB/s and 15~MB/s, respectively. The device operates within a supply voltage range of 2.7--3.6~V and an operating temperature range of 0\degree C to 70\degree C.

The card interfaces with the processing core unit via the SPI (Serial Peripheral Interface), using the MOSI, MISO, SCK, and CS signals for data transfer and control.

\subsection{Power Distribution Board}

%\begin{figure}[H]
%    \centering
%    \includegraphics[width=0.3\textwidth]{0_images/figures/Stepdown DCDC.png}
%    \caption{LM2596HVS-ADJ step-down DC/DC regulator}
%    \label{fig:component_LM2596HVS-ADJ_DC/DC_regulator}
%\end{figure}
Power distribution within the experiment is handled by a dedicated Power Distribution Board (PDB). The board receives a nominal input voltage of 28.8 V from the BEXUS power system and distributes regulated voltages to the experiment subsystems.

Voltage regulation is implemented using three LM2596HVS-ADJ step-down DC/DC converters. These regulators operate at a switching frequency of 150 kHz, support input voltages up to 57 V, and provide adjustable output voltages in the range of 1.2–37 V. The converters generate three regulated voltage rails required by the experiment: 24 V for the thermal and pressurization systems (compressor and vacuum pumps), 8 V for the K96 NDIR sensor and the MCU, and 3.3 V for the environmental sensors.

To mitigate switching noise introduced by the DC/DC converters on the shared power lines, an LC filter with a cut-off frequency of approximately 4 kHz is implemented on the power distribution board.

Due to the high startup currents expected from the pressurization system, a pre-charge circuit is implemented on the power distribution board to limit inrush currents during system startup.



\subsection{K96 Sensor}

\begin{figure}[H]
    \centering
    \includegraphics[width=0.5\textwidth]{0_images/figures/K96 Sensor.png}
    \caption{K96 Sensor. Taken from \cite{wastine2022}}
    \label{fig:component_K96_sensor}
\end{figure}

The K96 sensor is a \gls{ndir} gas sensor capable of measuring methane, carbon dioxide, and water vapor. It operates from a 6.0--8.5 V VBB supply and a 2.4--3.6 V IOVCC/RxD logic supply over a temperature range of $-40\degree$C to +85\degree C, with a maximum peak current of 300 mA. The sensor also provides local temperature, pressure, and humidity measurements via an integrated BME280 sensor.

\subsection{Temperature and Humidity Sensor}

%\begin{figure}[H]
%    \centering
%    \includegraphics[width=0.3\textwidth]{0_images/figures/Humidity Sensor.png}
%    \caption{SHT3X-DIS Humidity Sensor}
%    \label{fig:component_SHT3X-DIS_humidity}
%\end{figure}

The ambient temperature and humidity sensor is a SHT45, with an operating range of -40 to 125$^\circ$C and a humidity measurement range of 0--100\% RH. The sensor provides a humidity accuracy of $\pm 1.0$ \% RH and a temperature accuracy of $\pm 0.1\,^\circ$C. Communication with the sensor is performed via an I$^2$C digital interface, allowing temperature and humidity data to be read directly by the processing unit. The operating voltage range is 1.08 V to 3.6 V, with a typical average current consumption of approximately 0.4 $\mu$A.

\subsection{Pressure Sensor}
The ambient pressure sensor MS5607 is a barometric sensor that provides a highly linear
and precise 24-bit pressure and temperature value. It has I$^2$C and SPI interface up to
20MHz and has can measure between 10-1200mbar with an accuracy of ±1.5mbar. The
operating temperature range is -40 to 85\degree C and the supply voltage is between 1.8 and
3.6V

\subsection{Pressurization system}
The pressurization system is powered through the experiment power distribution board and controlled by the processing unit. The system includes two diaphragm vacuum pumps and a piston compressor, which together provide the airflow and pressure increase required for the measurement chamber. The system operates at 24 V and is supplied through the experiment power distribution system.

Pressure sensors will be installed between the compression stages and inside the measurement chamber to monitor system performance. These sensors provide feedback to the processing unit, enabling monitoring of the internal pressure and verification of correct system operation.

A passive outlet is used to allow air to exit the measurement chamber. The outlet is implemented as a fixed flow restriction that enables continuous exhaust of air while the pumps and compressor maintain the required chamber pressure. The internal chamber pressure is therefore primarily regulated by the pressurization system, while the passive outlet ensures continuous gas exchange within the chamber.

Detailed electrical interfaces, component specifications, and control parameters will be defined during the detailed design phase following component selection and testing.

%%%%%%%%%%%%%%%%%%%%%%%%%%%%%%%%%%%%%%%%%%%%%%%%%%%%%%%%%%%%%%%%%%%%%%%%%%%%%%%
\section{Thermal Design}\label{sec: thermal design}
\begin{figure}[H]
    \centering
    \includegraphics[width=0.8\textwidth]{0_images/figures/pressure_temperature_profile.png}
    \caption{Averaged pressure and temperature profile obtained by the Sodankylä weather station in Finland over all days of October 2025.}
    \label{fig:temp_press_profile}
\end{figure}

The temperature conditions in a stratospheric balloon flight are quite extreme with temperatures reaching as low as -80\degree C. At the same time the air pressure can drop as low as 10 mbar, which also implies that convective heat transfer is very weak, which might cause a poorly designed system to overheat. An averaged ambient pressure and temperature profile is shown in \fref{fig:temp_press_profile}. The low temperatures might cause substantial convective cooling at lower altitudes of around 8 km, while the low pressure in combination with solar radiation might cause substantial internal and radiative heating. 
In the following, the worst hot and cold cases will be analysed, and expected temperatures of the experiment will be shown by referencing the expected pump operations schedule.\par 

\subsection{Thermal analysis background}
A detailed description is given in \aref{app:thermal_background}.


\subsection{Cold case}\label{sec: thermal cold case}
The cold case is a steady state assumption, given environmental conditions at a specific altitude. The appropriate altitude is determined as the maximum of the air heat convection $Q_\mathrm{conv}(h)$ profile from the averaged temperature and pressure profiles given in \fref{fig:temp_press_profile}. It is assumed that no solar radiation is impacting the system (except \gls{ir} radiation from the Earth). Ambient pressure is taken from \fref{fig:temp_press_profile}, while the ambient temperature is set to be $-80$°C. Notably, the tubes containing pressurised, hot air, are assumed to be perfect thermal insulators, which is far from accurate. However, as the heat flux provided by the heated air is very small, this is acceptable.\par 
To estimate impact of different operations and system designs, the analysis takes into account whether pumps are turned on or off, whether the experiment has heaters or not and whether the experiment is insulated or not. The insulation is assumed to be a 5 cm thick box of styrofoam, completely encapsulating the chassis. The pumps are assumed to draw 50 W. This corresponds to the maximum power draw possible by two diaphragm pumps and one piston compressor. The active pumps were chosen according to the operations schedule as outlined in \sref{sec: pump system arch}. Assuming an efficiency of 60\% for the diaphragm pumps and 40\% for the piston compressor, the resulting heat load is 26 W.\par 

In \tref{tab:thermal results cold case} the results of the cold case analysis are displayed. 
Without the styrofoam insulation, the chassis reaches dangerously low temperatures, even with an active heating system. The heating system then would require over 75 W, which is unfeasible. Even with the styrofoam insulation, the required heating power is around 22 W, which is unsustainable for longer periods of time. However, when the pumps are being used, their heat alone is enough to keep the chassis and chamber within desired temperatures. The cold case analysis therefore concludes a necessity of insulation and active pumps.


\begin{table}[H]
\centering
\caption{Thermal Performance Analysis: Cold Case Scenarios at 8.8\,km Altitude}
\label{tab:thermal results cold case}
\begin{tabular}{@{}lll S[table-format=-2.2] S[table-format=-2.2] S[table-format=-2.2]@{}}
\toprule
\textbf{Insulation} & \textbf{Pumps} & \textbf{Thermal System} & {$T_{Chassis}$ [\si{\celsius}]} & {$T_{Chamber}$ [\si{\celsius}]} & {Power [\si{W}]} \\ \midrule
None & Off & None & -79.27 & -79.27 & {---} \\
None & Off & Active & -29.02 & 21.00 & 75.23 \\ 
Styrofoam & Off & Passive & -78.70 & -78.70 & {---} \\
Styrofoam & Off & Active & 6.52 & 21.00 & 21.79 \\ \midrule
None & Active & None & -61.80 & -61.80 & {---} \\
None & Active & Active & -20.40 & 21.00 & 62.28 \\ 
Styrofoam & Active & Passive & -72.41 & -72.41 & {---} \\
Styrofoam & Active & Active & 21.14 & 21.00 & -0.21 \\ \bottomrule
\addlinespace
\multicolumn{6}{l}{\small \textit{Note: Negative power values indicate active cooling required to maintain target temperature.}}
\end{tabular}
\end{table}

\subsection{Hot case}\label{sec: thermal hot case}
Similar to the cold case, the hot case will follow a lumped system approach. However, the appropriate altitude here is chosen to minimize the impact of heat convection via air and maximizing solar heating. This means that the hot case will assume environmental conditions like those found at 30 km altitude, namely $T=-66$°C, $p=0.01$ bar and cloud-free skies. As in the cold case analysis, the effects of pumps, active heating, and styrofoam insulation are tested for individually. The results are displayed in \tref{tab:thermal results hot case}. Without styrofoam insulation, and no thermal system/heating or active pumps, the chassis temperature will reach around $-3$°C. 
Adding the insulation leads to a counterintuitive temperature reduction to $-47$°C. This is due to the styrofoam reflecting more sunlight than the aluminium, while simultaneously loosing heat more efficiently via \gls{ir} radiation. 
As the heat loads involved are comparatively small, an active heating system has to provide around 14 W to raise the chamber temperature to 21°C. However, if pumps are active at 30 km altitude, the chassis will quickly reach temperatures higher than 21°C, requiring cooling of around 9 W instead.\par 

In conclusion, the hot case analysis shows that active pumps at high altitudes will necessitate a cooling system. In order to avoid complexity, it is therefore favourable to avoid using the pumps at high altitudes. 
This agrees well with the planned operations for the pumps, as given in \sref{sec: pump system arch}.

 


\begin{table}[H]
\centering
\caption{Thermal Performance Analysis: Hot Case Scenarios at 30\,km Altitude}
\label{tab:thermal results hot case}
\begin{tabular}{@{}lll S[table-format=-2.2] S[table-format=-2.2] S[table-format=-2.2]@{}}
\toprule
\textbf{Insulation} & \textbf{Pumps} & \textbf{Thermal System} & {$T_{Chassis}$ [\si{\celsius}]} & {$T_{Chamber}$ [\si{\celsius}]} & {Power [\si{W}]} \\ \midrule
None & Off & None & -3.16 & -3.16 & {---} \\
None & Off & Active & 16.39 & 21.00 & 6.93 \\ 
Styrofoam & Off & Passive & -47.32 & -47.32 & {---} \\
Styrofoam & Off & Active & 11.89 & 21.00 & 13.71 \\ \midrule
None & Active & None & 65.10 & 65.10 & {---} \\
None & Active & Active & 30.25 & 21.00 & -13.92 \\ 
Styrofoam & Active & Passive & -29.47 & -29.47 & {---} \\
Styrofoam & Active & Active & 26.75 & 21.00 & -8.65 \\ \bottomrule
\addlinespace
\multicolumn{6}{l}{\small \textit{Note: Negative power values indicate active cooling required to maintain desired chamber temperature.}}
\end{tabular}
\end{table}


\subsection{Continuous operation}
As before, the continous operation usses a lumped system approach. However, now the total energy of the system is not assumed to be constant, but transient. This means that the system temperature is changing with the changing environmental and internal heat loads. Therefore, thermal inertia needs to be taken into account via
\begin{align}
    mc_{p}\dot{T} =& Q_\mathrm{ele}+Q_\mathrm{pump}(h)+Q_\mathrm{transfer} + Q_\mathrm{rad\,in}(h) + Q_\mathrm{insulation} - Q_\mathrm{conv}(h) - Q_\mathrm{rad\,out} \label{eq: heat diff chassis}
\end{align}
for the chassis. As before, the power required to keep the pressure chamber at 21°C can be calculated via \eqref{eq:heat_balance_chamber}, since $\dot{T}$ is supposed to be equal to 0 for the chamber. 
In this scenario the pumps operate by the operating plan specified in \sref{sec: pump system arch}. Additionally, pumps and thermal systems are activated 30 minutes before launch to simulate pre-heating and pre-pressurisation of the chamber. Otherwise the pumps are assumed to work on lower power draw settings at altitudes under 14 km, adding an estimated heat load of 18 W. Between 14 km and 22 km, the pumps operate at full power, adding a heat load of 26 W. Beyond 22 km the pumps are switched off. The total flight time is estimated to be 6 h long. 
The float phase is assumed to be at constant 26 km altitude. The heating system is assumed to be capable of a maximum heating output of 10 W.
After the pumps are deactivated at 22 km, the heating system aims to stabilise the chamber temperature at $-5$°C. This is done as a proxy to prevent the electronics in the chassis from freezing.\par 

The result for an experiment inside a styrofoam insulation box and subject to solar daylight is shown in \fref{fig:continous thermal sim main}. 
%As expected, heating is especially necessary between ca 8 km and 14 km. 
The chassis temperature never drops much below $-10$°C, while the thermal power requirements do not reach the set maximum of 10 W.  
By integrating the power requirement curve one obtains a total required heating power of $\approx 29.45$ Wh. This result refers to heating energy and not to the energy necessary to power heating components, which is expected to be greater.\par 
If no insulation were used, the total required heating power would jump to 37.45 Wh. If insulation were used, but no daylight were available, the heating power required would be 36.42 Wh. The temperature and heating curves for both cases are provided in \fref{fig:continous thermal sim side}. If no insulation is used, the heating system is unable to keep the chamber temperature within the required temperature range of 15°C to 40°C during ascend. It also fails at keeping the chassis temperature within a safe range of -15 to 65°C during descend. On the other hand, if insulation is used, but the flight is during night time or particularly cloudy conditions, the effects are rather minimal. The heater would max out its power output and the chassis temperature would drop to around $-15$°C during descend, which is just the lower limit of the safe range and still 10°C away from the strict lower limit of -25°C.

\begin{figure}[htbp]
    \centering
    \includegraphics[width=0.8\textwidth]{0_images/figures/thermal_analysis_plot_WITHins_WITHlight.png}
    \caption{Expected temperature of the chassis and chamber and the heating power necessary to sustain 21°C in the chamber. The gray lines indicate the time intervals when the balloon is ascending, floating, and descending. The colors indicate different power settings of the pressurisation system: green - no pumps, yellow - low power, red - high power.}
    \label{fig:continous thermal sim main}
\end{figure}

\begin{figure}[htbp]
\begin{subfigure}[c]{0.5\linewidth}
\includegraphics[width=\linewidth]{0_images/figures/thermal_analysis_plot_NOins_WITHlight.png}
\caption{No insulation, with daylight.}
\end{subfigure}
\hfill
\begin{subfigure}[c]{0.5\linewidth}
\includegraphics[width=\linewidth]{0_images/figures/thermal_analysis_plot_WITHins_NOlight.png}
\caption{With insulation, no daylight}
\end{subfigure}%
\caption{Expected temperature of the chassis and chamber and the heating power necessary to sustain 21°C in the chamber. The gray lines indicate the time intervals when the balloon is ascending, floating, and descending. The colors indicate different power settings of the pressurisation system: green - no pumps, yellow - low power, red - high power.}
\label{fig:continous thermal sim side}
\end{figure}

\subsection{Conclusion}\label{sec: thermal conclusion}
The cold and hot case analysis shows that the experiment if enclosed in a styrofoam insulation is able to stay within the required temperature range depending on the proper use of pumps. In the ascend phase, the pumps are generating enough heat to keep the chamber temperature at 21°C, even without additional heating elements. In fact the transientt analysis shows that temperatures of chassis and chamber might reach close to 40°C. During float phase, the chamber temperature is allowed to drop to $-5$°C. This is done in order to keep the electrical components within the chassis at around $-10$°C. The required total heating power for daytime and nighttime would be $\approx 29.45$ Wh and $\approx 36.42$ Wh, respectively. Depending on the final design of the electrical components, the temperature limit and thus required power might be lowered further.

\subsection{Passive thermal control system}
The above analysis shows that white styrofoam is a good candidate for insulation. The chassis is planned to be completely sourrounded by 5 cm thick styrofoam insulation. Further analysis on other insulation materials and implementations, as well as on other passive control systems, will commence with the start of thesis project 6 and 7 (see \tref{tab:thesis}). Additionally, the intake air will be preheated before being fed to the diaphragm pumps. This can be done via a thermally conductive, solenoid shaped tubing housed inside the chassis. If the tubing has higher heat capacity than the incoming air, it can act as a heat sink, heating the incoming air to an acceptable temperature. The exact shape and length of the tube will be determined via tests.

\subsection{Active thermal control system}
The above analysis shows that during ascend and before launch, the pumps alone will provide enough heating power to heat the chamber to around 40°C. This is on the edge of the acceptable temperature range as outlined in \tref{tab:design_req}, P.7. Further analysis of overheating mitigation methods, including pump duty cycles and other active thermal control methods will commence with the start of thesis project 6 (see \tref{tab:thesis}).\par 

During float and descend, the above analysis shows that a 10 W heater is sufficient to keep the electronics inside the chassis within acceptable limits, see \tref{tab:design_req}, D.1.3. There are multiple common choices for thermal control components at low wattages and high precision. These will be discussed below.
\paragraph{\gls{ptc} ceramics} are commonly used in portable space heaters. They are used via resistive heating. However, due to their \gls{ptc}, a rising internal temperature causes their resistivity to increase as well. This behaviour can be configured, such that a desired maximum temperature can be set at which the ceramic resistivity becomes too high for the temperature to rise further. They are therefore self regulating, adding additional security and requiring less controlling. They are also quite cheap. However, because they are made of ceramic, they can not easily fit inside a tight experiment. 

\paragraph{Peltier devices} are heat pumps, that transfer heat from one side of the device to the other when an electrical current is applied. Upon reversal of the current polarity, the heat transfer flow is also reversed. These devices can therefore be used to cool and heat. This makes them great candidates for an active thermal control system that has to cool and heat components depending on mission phase. The downside is that they require heat sinks that the excess heat (or lack thereof) can be dumped to. They also require more sophisticated controlling software and circuits to be able to switch between cooling and heating.

\paragraph{Polyimide film heating elements} consist of a thin conductive heating element sourrounded by polyimide film. They are very flexible, sturdy and commonly used in space environment heating solutions. Although they are not self regulating, they can be easily controlled using \gls{pid} algorithms. They are unable to provide cooling, but offer a great surface area to enable uniform heating. Due to their flexibility, they are also easy to integrate into a system and do not require complex software or circuits.\par  

Because of the low heating power required, the small operating temperature expected, the ease of integration, its big surface area, and its low cost a polyimide film heating element is chosen. The heating element is included in \tref{tab:power_consumption}. It can be attached via \gls{psa} and wrapped around the pressure chamber. Since the chamber is made out of aluminium, and aluminium is good thermal conductor, one heating pad should suffice. This will be confirmed by tests. Should tests show that there is indeed no need for heating the pressurisation chamber, the heater may be moved directly to the temperature sensitive electrical components like the \gls{pcb}.\par 

The heating pad will be controlled using a \gls{pid} algorithm, as common for heating appliances. The algorithm will run on the \gls{mcu}. The \gls{mcu} then controls the heating power of the heating pad by adjusting its duty cycle by powering a \gls{mosfet}. The exact \gls{pwm} setting is to be determined by the final \gls{mcu} design. To prevent overheating of the chamber, a \gls{tco} unit is used, wired in series with the heating pad. It can be directly attached via \gls{psa} to the outside of the pressure chamber, monitoring its temperature and breaking the circuit when its temperature reaches above 50°C.


% This can lead to overheating of the electrical components if the heat from dissipated power accumulates. Therefore, it is essential to know how much power is dissipated and to take additional cooling measures in case that radiative cooling is insufficient.

% \begin{table}[H]
% \centering
% \caption{Component Temperature and Power Ratings}
% \label{tab:component_specs}
% \begin{tabular}{p{4cm}p{1cm}p{1cm}p{1cm}p{1cm}p{1cm}p{1cm}}
% \toprule
% \textbf{Component} &
% \multicolumn{2}{c}{\makecell[c]{Operating T (°C)}} &
% \multicolumn{2}{c}{\makecell[c]{Survivable T (°C)}} &
% \multicolumn{2}{c}{\makecell[c]{Power Dissipated (W)}} \\
% \cmidrule(lr){2-3} \cmidrule(lr){4-5} \cmidrule(lr){6-7}
% & Min & Max & Min & Max & Mean & Max \\
% \midrule
% Pump &  &  &  &  &  &    \\
% Valve &  &  &  &  &  &    \\
% Arduino Mega & -40 & 85 &  &  &  &  1.95  \\
% Switching Regulator & -40 & 125 &  &  &  & 39   \\
% SD Card & -25 & 85 &  &  &  &    \\
% K96 sensor (\gls{ndir}) & 0 & 40 & -40 & 85 &  & 2.447   \\
% BME sensor & -40 & 85 &  &  &  & 0.001   \\
% MS5607 (P) & -40 & 85 &  &  &  & 0.00004   \\
% DS18B20 (T) & -55 & 125 &  &  &  &    \\
% SHT3X-DIS (RH) & -40 & 125 &  &  &  &    \\
% \bottomrule
% \end{tabular}
% \end{table}

% The experiment can be thermally divided into three parts. The coldest part are the ambient sensors on the outside of the experiment box. They are and should be completely exposed and therefore are required to withstand temperatures as low as -80\degree C. This poses a challenge since most commercially available components are made for ground level conditions and manufacturers usually draw the line at -40\degree C for the temperature conditions. If we should get access to a thermal chamber, we can test the reliability of those sensors ourselves at lower temperatures. Since we haven't found sensors with operating ranges going down as low as -80\degree C we might have to resort to this attempt. Those sensors are not central to our experiment. The second part is the electronics compartment, which will likely be not actively heated. The bigger issue for this part might be overheating instead and a careful thermal analysis during the design process will be required to decide on mitigation steps if the temperature ranges of the components are exceeded at the upper or possibly still lower bounds. Later, in the integration process, further tests will follow to verify the functionality of that part at high altitudes.

% The last part is the measurement chamber of which the temperature needs to be more carefully regulated to ensure that the K96 \gls{ndir} sensor, which is central to our experiment, functions properly. Temperature fluctuations can also lead to output drift and therefore we need to keep the temperature in a narrow range. At the same time, pumps, valves or inlet and outlet tube openings could freeze over if the humidity level is high. Since we might fly through clouds with very high humidity levels it is essential that the relative humidity doesn't get too high. By increasing the temperature of the air, condensation can be mitigated and freezing of pumps, valves and tubes should be prevented as well. Further tests will be conducted in later stages of experiment integration.

% Insulation will play another essential role in the thermal design. The whole experiment should be as much thermally separated from the rest of the vehicle as possible. On one hand, this avoids further cooling of the experiment and on the other it avoids impacting other experiment or the vehicle. The measurement chamber should be well insulated as well to not (further) heat up the electronics compartment and to avoid condensation of water on the walls of the chamber.

%%%%%%%%%%%%%%%%%%%%%%%%%%%%%%%%%%%%%%%%%%%%%%%%%%%%%%%%%%%%%%%%%%%%%%%%%%%%%%%
\section{Power design}

The \gls{mirage} team expects to rely exclusively on the 28.8V SAFT MP176065 which is provided by the gondola.

The experiment will have an estimated average current of 0.80A and use 138.67Wh in total, see table 4.9, including the thermal system and the pressurization system. This is higher than previous estimations, mainly due to the thermal system, compressor and vacuum pumps, but this calculation is based on full capacity so it's expected to be lower than this. The peak current draw is high at 2.26A from the BEXUS side for ~30 minutes, see table 4.8, while the compressor and vacuum pumps are working at full capacity, which is close to the requirement D.11.1. 
Since the MCU PCB will be designed by thesis students it is difficult to estimate precisely the current draw for this component, however it was estimated to be 300mA. The K96 has a max current of 300mA but the expected average is lower depending on the time per measurement. 

Before launch, the pressurization system, e.g. vacuum pumps and compressor currents will be limited using some sort of pre-charge circuit and will thus not have a major impact on the energy consumption as a whole. During the flight phase, the pressurization system will be activated in lesser capacity for $\approx 45$ minutes and in full capacity $\approx 30$ minutes in order to conserve energy for the rest of the experiment. After the measurement phase the pressurization system will be limited or completely shut off.

The thermal system will not be activated while the pressurization system is active in order to not exceed current draw limitations and thus will only be operational as detailed in thermal design section.

\begin{table}[H]
\centering
\caption{Max current consumption}
\label{tab:current_consumption}
\begin{tabular}{llllll}
\toprule
\textbf{Amount} & \textbf{Components}& \textbf{Voltage} & \textbf{Current} &\textbf{Power} &\textbf{Energy}\\
\midrule
1 & MCU & 8V & 300mA & 2.4W & 1.2Wh\\
1 & K96 & 8V & 300mA & 2.4 & 1.2Wh \\
2 & MS5607 & 3.3V& 1µA & 6.6µW & 3.3µWh \\
6 & SHT45 & 3.3V & 0.4µA & 7.92µW & 3.96µWh \\
1 & Compressor & 24V & 1.25A & 30W & 15Wh \\
2 & Vacuum pumps & 24V & 0.42A & 20W & 10Wh\\
1 & SDHC & 3.3V & 150 mA & 0.495W & 0.248Wh \\
\midrule
Total &&& 2.45A & 55.295W & 27.648Wh\\
Total (BEXUS) &&28.8V& 2.26A & 65.05W & 32.526Wh \\
\bottomrule
\end{tabular}
\end{table}

\begin{table}[H]
\centering
\caption{Power consumption from the power distribution board and the Bexus power unit}
\label{tab:power_consumption}
\begin{tabular}{L{1.6cm}llllll}
\toprule
\textbf{Duration} & \textbf{Amount} & \textbf{Component} & \textbf{Voltage} & \textbf{Current} & \textbf{Power} & \textbf{Total Energy} \\
\midrule
6h & 1 & MCU & 8V & 300mA & 2.4 W & 14.4Wh \\
1.25h & 1 & K96 & 8V & 300mA & 2.4 W & 3Wh \\
6h & 6 & SHT45 & 3.3 V & 0.4µA & 1.32 µW & 47.52µWh \\
1.25h & 1 & Compressor & 24V & 1.25A & 30W & 37.5Wh \\
1.25h & 2 & Vacuum Pump & 24V & 0.42A & 10W & 25Wh \\
6h & 2 & Pressure sensor & 3.3V & 1µA & 3.3µW & 39.6µWh \\
6h & 1 &SDHC &3.3V& 150mA & 0.495W & 2.97Wh \\
3.5h & 1 &Heating element &12V& 0.83A & 10W & 35Wh \\

\midrule
Average (Bexus) & & & 28.8V & 0.80A & 23.11W\\
Total Bexus Power & & & & & & 138.67Wh \\
\bottomrule
\end{tabular}
\end{table}

Table 4.9 presents key components, amount of components, duration of use, voltage levels, currents and power. The table is representing from the experimenters perspective while the average and total BEXUS power is from the BEXUS power unit perspective. This calculations used the assumption of $85\%$ efficiency to calculate the power consumtion on the BEXUS power unit. 
%%%%%%%%%%%%%%%%%%%%%%%%%%%%%%%%%%%%%%%%%%%%%%%%%%%%%%%%%%%%%%%%%%%%%%%%%%%%%%%
\section{Software design}

The on-board software is responsible for collecting data from the sensors and store it on the on-board SD card. It is also responsible for handling communication from and to ground, as well as the control of the internal pressure and thermal systems and preventing condensation.

For the pressurisation system the software needs to be able to read out the pressure sensor in the measurement chamber and adjust the power of the pumps accordingly to maintain a constant pressure of 3 bar. The mechanical side of the system is explained above. Since the vacuum pumps prepressurise the air that is then compressed to 3 bar, the power of the vacuum pumps and the compressor need to be handled separately. Depending on whether the compressor needs to have a minimum flowrate to work, it might be possible that a pulsed operation is necessary for the compressor, which would put strain on the electrical system. The software needs to be able to adjust these different components accordingly based on the internal pressure sensor. \\

For the thermal control system the software needs to take sensor inputs from different places. It should be split into two different programs, where one is responsible for stabilizing the temperature in the measurement chamber at around room temperature. This also includes controlling potential active heating elements at the gas inlet of the outer enclosure and before the air at 3 bar enters the measurement chamber based on at least on the thermal sensor of the K96 but possibly also several thermal sensor in between. The other thermal control loop should monitor components that are at risk of overheating and components that could underheat. There might be also active heating elements at critical components like the MCU, which could be switched on if the temperatures drop dangerously low inside the electrical compartment. This control loop need to take the different temperature outputs of sensors that are close to the pumps and electrical components and if necessary take counter measures via controlling the heating elements or reducing the power of the pumps.\\
For humidity the role of the software is to detect to high humidity levels (>80\% r.H.). If the ambient sensors detect high levels of humidity outside the software should close the shutter and stop the pumps as a first step to prevent more air with high humidity levels from entering the measurement chamber without losing pressure. When the humidity levels are low enough, the temperature of the air taken into the measurement chamber should be increased slightly by increasing the temperature of the active heating elements until the internal humidity is lower as well.\\


The settings of those systems depend on the readings of the internal and external pressure and temperature sensors. Additionally, correct flight phases need to be engaged. This is done based on longer term changes in the recorded ambient pressure or via a command from the ground, with the command taking priority. A control flow and state diagram are provided in Figure \ref{fig:cont_flow} and Figure \ref{fig:state_diagram}.

The experiment will run on 2 main modes: Manual and Autonomous. Each of the main modes has a submode for standby and active measurements. When the experiment reaches the point where no more data needs to be collected, the system will move to standby mode unless a command has been given to do otherwise. The modes will also allow for autonomous work even through radio blackouts and loss of communication. If below 22\,km and a loss of connection occurs, the system will automatically switch to autonomous measurements mode; if however the experiment is above 22\,km and a loss of connection occurs then the system will switch to Autonomous standby mode that will ensure the experiment remains safe without exhausting any of the hardware.

To guarantee operations watchdog systems will be used to reset the operations if it encountered a fatal error.
The software system will work with different threads working asynchronously.
Data transfer between the threads will happen through producer-consumer queues that are thread-safe.
There will be a thread for reading sensor data, a thread of controlling hardware, a thread sending downlink packages and receiving packages from uplink, a thread to write data on the SD card and finally a thread to control everything with the aid of a watchdog.

All transmissions are saved to the on-board storage with timestamps. Similarly, all saved data is saved with timestamps, too. A table detailing the expected data sizes of data to be stored is given in Table \ref{tab:data_rates}. In total 39856 bits per second (40 kbit) will be stored. This means that for the duration of 6 hours a storage of 108 MB (megabyte) is required to store all the data received from the sensors. The SD cards onboard are more than enough for this task. Some of this data will be downlinked every seconds for redundancy and tracking of the functionality of the experiment. Based on the sampling rates in Table \ref{tab:data_rates},the Elink thread can pull 100 datapoints from the telemetry and data queue to downlink each second, resulting in an estimated downlink rate of 32.8 kbps, below the experiment requirements (P.18) of 100 kbps. Assuming a maximum connection dropout duration of 10 minutes, this would likewise imply a build-up of data of 19.68 Mbit. Transferring it over the next 5 minutes gives a data rate of 65.6 kbps which is feasible. With compression it is probable that only half the rates stated need to be transmitted, but a higher data rate could be used to send duplicates or checksums for error mitigation.

The frequencies are chosen based on the datasheets of the sensors. The SHT45 sensor has multiple modes, fast mode plus (FM+) was chosen for this mission. The K96 sensor relies on the speed of the NDIR sensor within with typically is below 10 Hz, whereas in the MS5607 sensor OSR 4096 mode was chosen. The \gls{rtc} frequency was calculated to be equal to all the frequencies added. This is to consider the scenario where none of the sensors data readings collide and all the readings from the 3 sensors are collected at different timestamps, which will require one timestamp for each.

The K96 sensor core requires a UART interface and relies on the Modbus protocol, while most other sensors use I2C. The communication protocol will be UDP.

\begin{table}[H]
\centering
\caption{Max sampling rates, sampling sizes, and data rates of data products to be stored}
\label{tab:data_rates}
\begin{tabular}{p{7cm}p{2cm}p{2cm}p{2cm}}
\toprule
\textbf{Component} & \textbf{Max. sampling rate [Hz]} & \textbf{Sampling size [bit], includes CRC} & \textbf{Data rate [bps]} \\
\midrule
LPL (K96) & 10 & 32 & 320\\
SPL (K96) & 10 & 32 & 320\\
MPL (K96) & 10 & 32 & 320\\
Pressure (K96) & 10 & 16 & 160\\
Temperature (K96) & 10 & 16 & 160\\
Relative Humidity (K96) & 10 & 16 & 160\\
Internal pressure sensor (MS5607) & 110 & 28 & 3080 \\
Internal temperature sensor (SHT45) & 222 & 24 & 5328 \\
Internal humidity sensor (SHT45) & 222 & 24 & 5328 \\
Ambient pressure sensor (MS5607) & 110 & 28 & 3080 \\
Ambient temperature sensor (SHT45) & 222 & 24 & 5328 \\
Ambient humidity sensor (SHT45) & 222 & 24 & 5328 \\
Real Time Clock & 342 & 32 & 10944\\
% \dos{General housekeeping data} (current for e.g. pumps, data from thermal system, ...) & 1 & 384 & 384 \\
\midrule
Total & & & 39856 \\
\bottomrule
\end{tabular}
\end{table}

\begin{figure}[H]
\centering
\includegraphics[width=.9\textwidth]{0_images/figures/control_flow.png}
\caption{Control flow diagram}
\label{fig:cont_flow}
\end{figure}

\begin{figure}[H]
\centering
\includegraphics[width=\textwidth]{0_images/figures/state_chart.png}
\caption{State diagram}
\label{fig:state_diagram}
\end{figure}

%%%%%%%%%%%%%%%%%%%%%%%%%%%%%%%%%%%%%%%%%%%%%%%%%%%%%%%%%%%%%%%%%%%%%%%%%%%%%%%
\section{Ground support equipment}

A \gls{gui} will be developed for easy control of \gls{mirage} during the mission. The \gls{gui} will allow ground operators to send commands through uplink and visualize the downlink data. Additionally the team will require basic electronic equipment, chairs and a table to operate from and troubleshoot any issues before flight.

%%%%%%%%%%%%%%%%%%%%%%%%%%%%%%%%%%%%%%%%%%%%%%%%%%%%%%%%%%%%%%%%%%%%%%%%%%%%%%%\newpage

\chapter{EXPERIMENT VERIFICATION AND TESTING}

\section{Verification Matrix}

\begin{longtable}{L{0.7cm}L{5cm}L{1.5cm}L{2cm}L{2.5cm}L{2cm}}
\caption{Verification Matrix}
\label{tab:verification_matrix} \\

\toprule
\textbf{ID} & \textbf{Requirement text} & \textbf{Method} & \textbf{Reference} & \textbf{Status} & \textbf{Verification Result} \\
\midrule
\endfirsthead

\multicolumn{6}{c}{\tablename\ \thetable\ -- \textit{Continued from previous page}}\\
\toprule
\textbf{ID} & \textbf{Requirement text} & \textbf{Method} & \textbf{Reference} & \textbf{Status} & \textbf{Verification Result} \\
\midrule
\endhead

\midrule
\multicolumn{6}{r}{\textit{Continued on next page}}\\
\endfoot

\bottomrule
\endlastfoot

F.1. & The experiment shall measure the atmospheric CH\textsubscript{4}, CO\textsubscript{2} and H\textsubscript{2}O concentrations of ambient air simultaneously with an off-the-shelf \gls{ndir} sensor. & R, T & Test 3, Test 13 & To be done (June) & Pending \\
F.2. & The measurement chamber shall be pressurised with ambient air. & R, T & Test 2 & To be done (June) & Pending \\
F.3. & The air in the measurement chamber shall be exchanged. & R,T & Test 2, Test 3 & To be done (earliest test in June) & Pending \\
F.4. & The experiment shall measure pressure inside of the measurement chamber. & R,T & Test 2, Test 7 & To be done (earliest test in June) & Pending \\
F.5. & The experiment shall measure the temperature inside of the measurement chamber. & R,T & Test 3, Test 8  & To be done (earliest test in June) & Pending \\
F.6. & The experiment shall stabilize the temperature inside the measurement chamber. & A,T & Test 5, Test 8 & To be done (earliest test in June) & Pending \\
F.7. & The experiment shall measure the humidity inside of the measurement chamber. & R,T & Test 3, Test 8 & To be done (earliest test in June) & Pending \\
F.8 & The experiment shall prevent the condensation of water inside the measurement chamber. & R, T & Test 8 & To be done (July) & Pending \\
F.9. & The experiment should measure ambient temperature outside. & R,T & Test 8, Test 13 & To be done (earliest in July) & Pending \\
F.10. & The experiment should measure ambient pressure outside. & R,T &  Test 7, Test 13 & To be done (earliest in July) & Pending \\
F.11. & The experiment should measure ambient humidity outside. & R,T &  Test 8, Test 13 & To be done (earliest in July) & Pending \\
F.12. & The experiment shall send data via E-Link. & T & Test 6 & To be done (June) & Pending \\
F.13. & The experiment shall receive data via E-Link. & T & Test 6 & To be done (June) & Pending \\
P.1.1 & The \gls{ndir} sensor shall measure the concentration of methane with a range of 0.4 to 3.0 ppm. & R, T & Test 3, Test 13 & To be done (June) & Pending \\
P.1.2 & The \gls{ndir} sensor shall measure the concentration of methane with an uncertainty of 0.5 ppm or less. & R, T & Test 3, Test 13 & To be done (June) & Pending\\
P.2.1 & The \gls{ndir} sensor shall measure the concentration of carbondioxide with a range of 200 to 500 ppm. & R, T & Test 3, Test 13  & To be done (June) & Pending \\
P.2.2 & The \gls{ndir} sensor shall measure the concentration of carbondioxide with an uncertainty of 5 ppm or less. & R, T & Test 3, Test 13 & To be done (June) & Pending\\
P.3.1 & The \gls{ndir} sensor shall measure the concentration of water with a range of 1 to 30,000 ppm. & R, T & Test 3, Test 13  & To be done (June) & Pending \\
P.3.2 & The \gls{ndir} sensor shall measure the concentration of water with an uncertainty of 5 ppm or less. & R, T & Test 3, Test 13  & To be done (June) & Pending \\
P.4.1 & The pressure inside of the measurement chamber shall stay within a range of 2.5 and 3.2 bar.  & T & Test 2, Test 7 & To be done (earliest test in June) & Pending \\
P.4.2 & The pressure inside of the measurement chamber should stay within a range of 2.8 and 3 bar. & T & Test 2, Test 7 & To be done (earliest test in June) & Pending \\
P.5 & The pressurisation system shall have an air throughput of at least 1 l/min. & R, T & Test 2 & To be done (earliest test in June) & Pending \\
P.6.1 & The experiment shall measure the pressure inside of the measurement chamber in a range of 0.3 to 4 bar. & R, T & Test 2, Test 7 & To be done (June) & Pending \\
P.6.2 & The experiment shall measure the pressure inside of the measurement chamber with a maximum uncertainty of 0.02 bar. & R, T & Test 2, Test 7 & To be done (June) & Pending \\
P.7 & The temperature of the air inside of the measurement chamber shall stay within a range of 15 and 40\degree C while active. & A, T & Test 5, Test 8 & To be done (June)& Pending\\
P.8 & The temperature of the air inside of the measurement chamber should stay within a range of 19 and 24\degree C while active. & A, T & Test 5, Test 8 & To be done (June)& Pending \\
P.9.1 & The experiment shall measure the temperature inside of the measurement chamber in a range of -40 to +85 \degree C.  & R, T & Test 5, Test 8 & To be done (earliest test in June) & Pending \\
P.9.2 & The experiment shall measure the temperature inside of the measurement chamber with a maximum uncertainty of 1.5\degree C. & R, T & Test 5, Test 8 & To be done (earliest test in June) & Pending \\
P.10 & The experiment shall keep the humidity inside of the measurement chamber below 90\% r.H. & T & Test 8 & To be done (earliest test in July) & Pending \\
P.11.1 & The experiment shall measure the humidity inside of the measurement chamber in a range of 0 to 100\% r.H. & R, T & Test 3, Test 8 & To be done (earliest test in June) & Pending \\
P.11.2 & The experiment shall measure the humidity inside of the measurement chamber with a maximum uncertainty of 3\% r.H. & R, T & Test 3, Test 8 & To be done (earliest test in June) & Pending \\

P.12 & The experiment should measure the temperature, pressure and humidity inside of the measurement chamber at a rate over 10 Hz. & R, T & Test 8, Test 9 & To be done (earliest test in July) & Pending \\
P.13.1 & The experiment should measure the ambient pressure outside in a range of 20 to 1100 mbar. & R, T & Test 13 & To be done (August) & Pending \\
P.13.2 & The experiment should measure the ambient pressure outside with a maximum uncertainty of 20 mbar. & R, T & Test 13 & To be done (August) & Pending  \\
P.14.1 & The experiment should measure the ambient temperature outside in a range of -70 to +60 \degree C. & R, T & Test 13, Test 8 & To be done (earliest test in July) & Pending \\
P.14.2 & The experiment should measure the ambient temperature outside with a maximum uncertainty of 1.5\degree C. & R, T & Test 13, Test 8 & To be done (earliest test in July) & Pending \\
P.15.1 & The experiment should measure the ambient humidity outside in a range of 0 to 100\% r.H. & R, T & Test 13, Test 8 & To be done (earliest test in July) & Pending \\
P.15.2 & The experiment should measure the ambient humidity outside with a maximum uncertainty of 3\% r.H. & R, T & Test 13, Test 8 & To be done (earliest test in July) & Pending \\
P.16 & The experiment should measure the ambient temperature, pressure and humidity outside at a rate over 10 Hz. & R, T & Test 13, Test 9 & To be done (earliest test in July) & Pending \\
P.17 & The experiment should receive data via E-Link with an uplink rate under 10 kbps. & R, T & Test 12, Test 14 & To be done (July) & Pending \\
P.18 & The experiment shall send the obtained data with the E-Link with a downlink rate under 100 kbps. & R, T & Test 12, Test 14 & To be done (July)& Pending \\

D.1.1 & The experiment shall operate in the temperature profile of the BEXUS vehicle flight and launch. & A, T & Test 8 & To be done (July) & Pending \\
D.1.2 & The temperature in the electronic compartment shall stay within a range of -25 and 85°C. & A, T & Test 8 & To be done (July) & Pending\\
D.1.3 & The temperature in the electronic compartment should stay within a range of -15 and 60°C.& A, T & Test 8 & To be done (July) & Pending\\
D.2.1 & The experiment shall operate in the pressure profile of the BEXUS vehicle flight and launch. & A, T & Test 1, Test 2, Test 4, Test 7 & Test 1 completed successfully  & Pending \\
D.2.2 & The measurement chamber shall withstand a pressure difference of up to 4 bar.  & A, T & Test 7 & To be done (July) & Pending \\
D.3 & The experiment shall operate in the vibration profile of the BEXUS vehicle flight and launch. & T & Test 10 & To be done (July) & Pending \\
D.4 & The experiment shall operate in the shock profile of the BEXUS vehicle flight and launch. & T & Test 11 & To be done (July) & Pending \\
D.5 & The experiment power connector shall be an Amphenol PT02E8-4P, Code A. & R & & & \\
D.6 & The experiment data connector shall be an Amphenol RJF21B connector.  & R & & & \\
D.7 & The experiment shall not disturb or harm the launch vehicle.  & R, A &  &  & Pending \\
D.8 & Only the measurement chamber containing the \gls{ndir} sensor core and an additional pressure sensor shall be pressurised. & I, R & & & \\
D.9 & The experiment shall consume less than 150 Wh during the flight. & T, A & Test 7, Test 8, Test 14 & To be done (earliest test in July) & Pending \\
D.10 & The experiment shall operate under 28.8\,V.  & R, T & Test 7, Test 8, Test 14 & To be done (earliest test in July) & Pending \\
D.11.1 & The experiment shall have no currents exceeding 2.4 A. & R, T & Test 7, Test 8, Test 14 & To be done (earliest test in July) & Pending \\
D.11.2 & The experiment should have no currents exceeding 1.4 A.  & R, T & Test 7, Test 8, Test 14 & To be done (earliest test in July) & Pending \\
D.12 & The experiment should stay within a size limit of 15x20x20 cm\textsuperscript{3}. & T & Test 9 & To be done (July) & Pending  \\
D.13 & The experiment shall stay within a weight limit of 5 kg. & A,T & Test 9 & To be done (July) & Pending \\
D.14 & The air sample to be measured shall not be changed in its composition by the experiment, the vehicle or other experiments. & R, T & Test 13 & To be done (August) & Pending \\

D.15 & The experiment shall have access to fresh air from outside. & R, I & & & \\
D.16 & The data storage shall be able to store at least 2GB. & R & & & \\
D.17 & The data storage shall be protected from water due to condensation or during landing. & R & & & \\
D.18 & The \gls{ndir} sensor should still function after recovery post flight. & T & Test 9, Test 10 & To be done (earliest test in July) & Pending \\
D.19 & The budget for all components except the K96 sensor should stay below 2000\texteuro.  & R & & & \\
O.1 & The experiment shall be able to conduct measurements autonomously in case connection with the ground segment is lost. & T & Test 12, Test 14 & To be done (earliest test in July) & Pending \\
O.2 & The experiment shall accept a request for radio silence at any time while on the launch pad.  & T & Test 12, Test 14 & To be done (earliest test in July) & Pending \\
O.3 & The experiment shall store all data on board and continuously send data to ground via E-Link.  & T & Test 12 & To be done (July) & Pending \\
O.4 & The experiment shall automatically register and respond to high r.H. inside the chamber in order to lower the humidity. & T & Test 8, Test 13 & To be done (earliest test in July) & Pending \\

\end{longtable}

\section{Verification Plan}

%\begin{table}[htbp]
%\centering
%\caption{Controlled functionality testing and calibration of K96 sensor %only}
%\label{tab:test1}
%\begin{tabular}{L{5cm}L{10cm}}
%\toprule
%Test number & 1 \\
%Test type & Sensor functionality \\
%Test facility & Uppsala University \\
%Tested item & K96 sensor \\
%Model & Prototype Model \\
%Procedure, Test level and duration & Performance test, compare sensor %readings to known composition of gases tested with varying humidity %levels for measurement times of 20, 50 and 100s. \\
%Test campaign duration & 2 days \\
%Test campaign date & March \\
%Test completed & No \\
%Requirements verified & \\
%\bottomrule
%\end{tabular}
%\end{table}
\begin{table}[H]
\centering
\caption{Pressure test light bulb}
\label{tab:test_pressure_light_bulb}
\begin{tabular}{L{5cm}L{10cm}}
\toprule
Test number & 1 \\
Test type & Pressure test \\
Test facility & Uppsala University \\
Tested item & Light bulb of K96 sensor \\
Model & Prototype model \\
Procedure, Test level and duration & First part: Put Light bulb into vacuum chamber and drop pressure slowly to 10 mbar. \\
& Second part: Put light bulb into airtight container and increase pressure slowly until at least 4 bar are reached or until the light bulb bursts. \\
Test campaign duration & 2 days \\
Test campaign date & March \\
Test completed & Partially \\
Requirements verified & High pressure test sucessfully conducted\\
\bottomrule
\end{tabular}
\end{table}

\begin{table}[H]
\centering
\caption{Test of pressurisation system}
\label{tab:test_pressurisation_system}
\begin{tabular}{L{5cm}L{10cm}}
\toprule
Test number & 2 \\
Test type & Pressure test \\
Test facility & Uppsala University, Vacuum chamber \\
Tested item & Measurement chamber without K96\\
Model & Prototype model \\
Procedure, Test level and duration & Slowly lower pressure from 1 bar to 10 mbar and observe stability of pressure inside of the box. Simulate balloon ascent in duration and rate of pressure decline.  \\
Test campaign duration & 1 day \\
Test campaign date & June \\
Test completed & No \\
Requirements verified & \\
\bottomrule
\end{tabular}
\end{table}


\begin{table}[H]
\centering
\caption{Measurement chamber functionality}
\label{tab:test3}
\begin{tabular}{L{5cm}L{10cm}}
\toprule
Test number & 3 \\
Test type & Sensor functionality \\
Test facility & ICOS, Norunda \\
Tested item & Measurement Chamber \\
Model & Prototype Model \\
Procedure, Test level and duration & Put measurement chamber close to atmospheric sensing facilities at the location and cross reference collected data. Compare ambient sensor readings with measurements done by ICOS as well. \\
Test campaign duration & 2 days (setup and long term data collection) \\
Test campaign date & June \\
Test completed & No \\
Requirements verified & \\
\bottomrule
\end{tabular}
\end{table}


\begin{table}[H]
\centering
\caption{Vacuum test of the electronic compartment}
\label{tab:test4}
\begin{tabular}{L{5cm}L{10cm}}
\toprule
Test number & 4 \\
Test type & Pressure test \\
Test facility & Uppsala University, Vacuum chamber \\
Tested item & Electronic compartment \\
Model & Prototype model \\
Procedure, Test level and duration & Slowly lower pressure from 1 bar to 10 mbar and observe stability of circuit and temperature inside and outside of box. Simulate balloon ascent in duration and rate of pressure decline. Keep at 10 mbar for 2 hours at least. \\
Test campaign duration & 1 day \\
Test campaign date & June \\
Test completed & No \\
Requirements verified & \\
\bottomrule
\end{tabular}
\end{table}


\begin{table}[H]
\centering
\caption{Thermal control test}
\label{tab:test_thermal}
\begin{tabular}{L{5cm}L{10cm}}
\toprule
Test number & 5 \\
Test type & Thermal \\
Test facility & Thermal chamber and vacuum chamber \\
Tested item & Measurement chamber without K96 sensor \\
Model & Prototype model \\
Procedure, Test level and duration & In Thermal chamber: maintain room temperature of air in measurement chamber while decreasing temperature down to -40°C by autonomous current regulation. \\
& In Vacuum chamber: maintain room temperature of air in measurement chamber while decreasing pressure down to 10\,mbar by autonomous current regulation.  \\
Test campaign duration & 2 days \\
Test campaign date & June \\
Test completed & No \\
Requirements verified & \\
\bottomrule
\end{tabular}
\end{table}




\begin{table}[H]
\centering
\caption{Static loads test}
\label{tab:test_static_loads}
\begin{tabular}{L{5cm}L{10cm}}
\toprule
Test number & 6 \\
Test type & Mechanical test \\
Test facility & Uppsala University \\
Tested item & Main structure of experiment \\
Model & Qualification model \\
Procedure, Test level and duration & Mounting of experiment should simulate the gondola connection. Then the experiment is loaded with 10 to 30 times the experiment's own weight. \\
Test campaign duration & 1 day \\
Test campaign date & June (latest July) \\
Test completed & No \\
Requirements verified & \\
\bottomrule
\end{tabular}
\end{table}



\begin{table}[H]
\centering
\caption{Vacuum test of whole experiment}
\label{tab:vacuum_test}
\begin{tabular}{L{5cm}L{10cm}}
\toprule
Test number & 7  \\
Test type & Pressure test \\
Test facility & Uppsala University, Vacuum chamber \\
Tested item & Whole experiment \\
Model & Protoflight model \\
Procedure, Test level and duration & Slowly lower pressure from 1 bar to 10 mbar and observe stability of circuit and temperature inside and outside of box. Simulate ascent in duration and rate of pressure decline. Keep at 10 mbar for at least 15 minutes. Pressurise measurement chamber during the whole time. \\
Test campaign duration & 2 days (setup and execution) \\
Test campaign date & July \\
Test completed & No \\
Requirements verified & \\
\bottomrule
\end{tabular}
\end{table}


\begin{table}[H]
\centering
\caption{Temperature and humidity test of whole experiment}
\label{tab:test_temperature_humidity}
\begin{tabular}{L{5cm}L{10cm}}
\toprule
Test number & 8 \\
Test type & Thermal and humidity \\
Test facility & \gls{irf}, Thermal Chamber \\
Tested item & Whole experiment \\
Model & Proto flight model \\
Procedure, Test level and duration & Vary temperature and test functionality, 20 minutes at temperature steps of 5\degree C between -40\degree C and 10\degree C. Vary humidity between 5 and 100\% r.H. at each temperature step. \\
Test campaign duration & 3 days (1 day build-up, 2 day testing and packing) \\
Test campaign date & July \\
Test completed & No \\
Requirements verified & \\
\bottomrule
\end{tabular}
\end{table}


\begin{table}[H]
\centering
\caption{Weight and dimensions test}
\label{tab:weight_test}
\begin{tabular}{L{5cm}L{10cm}}
\toprule
Test number & 9 \\
Test type & Mechanical \\
Test facility & Uppsala University \\
Tested item & Whole experiment \\
Model & Qualification model \\
Procedure, Test level and duration & Put the experiment on a scale. Take a measurement tape and make sure that the outside dimensions are as psecified in the design. \\
Test campaign duration & 10 minutes \\
Test campaign date & July \\
Test completed & No \\
Requirements verified & \\
\bottomrule
\end{tabular}
\end{table}


\begin{table}[H]
\centering
\caption{Vibration test}
\label{tab:test_vibration}
\begin{tabular}{L{5cm}L{10cm}}
\toprule
Test number & 10 \\
Test type & Mechanical test \\
Test facility & Outside \\
Tested item & Whole experiment \\
Model & Qualification model \\
Procedure, Test level and duration & Mount experiment in the back of a car and drive through rough terrain, inspect condition of experiment afterwards and test full functionality of experiment. \\
Test campaign duration & 2 days \\
Test campaign date & July \\
Test completed & No \\
Requirements verified & \\
\bottomrule
\end{tabular}
\end{table}


\begin{table}[H]
\centering
\caption{Schock test}
\label{tab:test_shock}
\begin{tabular}{L{5cm}L{10cm}}
\toprule
Test number & 11 \\
Test type & Mechanical \\
Test facility & Uppsala University \\
Tested item & Whole experiment \\
Model & Qualification model \\
Procedure, Test level and duration & Drop experiment from 1 to 3 meters onto a padded surface and check visually and for full functionality afterwards. \\
Test campaign duration & 1 day \\
Test campaign date & July \\
Test completed & No \\
Requirements verified & \\
\bottomrule
\end{tabular}
\end{table}


\begin{table}[H]
\centering
\caption{Software testing E-Link}
\label{tab:test6}
\begin{tabular}{L{5cm}L{10cm}}
\toprule
Test number & 12 \\
Test type & Software test \\
Test facility & Uppsala University \\
Tested item & E-link, EGSE and software, measurement chamber \\
Model & Whole experiment \\
Procedure, Test level and duration & Send and receive data via E-link, test manual control of pumps, thermal control and measurement time \\
Test campaign duration & 1 day \\
Test campaign date & July \\
Test completed & No \\
Requirements verified & \\
\bottomrule
\end{tabular}
\end{table}


\begin{table}[H]
\centering
\caption{Functionality test with reference measurement}
\label{tab:test13}
\begin{tabular}{L{5cm}L{10cm}}
\toprule
Test number & 13 \\
Test type & Full referenced functionality test \\
Test facility & ICOS, Norunda \\
Tested item & Whole Experiment \\
Model & Flight Model \\
Procedure, Test level and duration & Put measurement chamber close to atmospheric sensing facilities at the location and cross reference collected data. Compare ambient sensor readings with measurements done by ICOS as well. \\
Test campaign duration & 2 days (setup and long term data collection) \\
Test campaign date & August \\
Test completed & No \\
Requirements verified & \\
\bottomrule
\end{tabular}
\end{table}

\begin{table}[H]
\centering
\caption{Bench test}
\label{tab:test_bench}
\begin{tabular}{L{5cm}L{10cm}}
\toprule
Test number & 14 \\
Test type & Long term test \\
Test facility & Uppsala University \\
Tested item & Whole experiment \\
Model & Flight Model \\
Procedure, Test level and duration & Assemble as for flight and run through simulated countdown of about 2 hours. Run experiment for 6 hours simulating ascent, float and descent. Simulate E-link dropouts as well. Consider recovery time as well for at least 24 hours. \\
Test campaign duration & 2 days \\
Test campaign date & August \\
Test completed & No \\
Requirements verified & \\
\bottomrule
\end{tabular}
\end{table}


% % ... Additional test tables would follow similarly ...

% \section{Verification Results}
% Test will be conducted later on and the results will be listed here in later versions of the SED.
% %
% %\begin{table}[htbp]
% %\centering
% %\caption{Example of a test result overview}
% %\label{tab:test_result}
% %\begin{tabular}{L{5cm}L{10cm}}
% %\toprule
% %\textbf{Verification number} & \\
% %\textbf{Type of test} & \\
% %\textbf{Facility} & \\
% %\textbf{Verified item} & \\
% %\textbf{Verification description} & \\
% %\textbf{Expected results} & \\
% %\textbf{Obtained Results} & \\
% %\textbf{Conclusions} & \\
% %\bottomrule
% %\end{tabular}
% %\end{table}
% %\newpage

\chapter{LAUNCH CAMPAIGN PREPARATIONS}
\label{sec:campaign_preparations}
\section{Input for the Campaign / Flight Requirement Plan}
\label{sec:input_FRP}

%%%%%%%%%%%%%%%%%%%%%%%%%%%%%%%%%%%%%%%%%%%%%%%%%%%%%%%%%%%%%%%%%%%%%%%%%%%%%%%
\subsection{Dimensions and Mass}

\begin{table}[htbp]
\centering
\caption{Experiment Mass and Volume}
\label{tab:mass_volume}
\begin{tabular}{L{6cm}L{6cm}}
\toprule
Experiment mass (in kg): & $\sim$4\\
Experiment footprint area (for BEXUS) (in $m^2$): & 0.04\\
Experiment volume (in $m^3$): & 0.006\\
Experiment expected \gls{cog} position: & Distance from geometric centre: x=4mm y=23mm z=22mm \\
\bottomrule
\end{tabular}
\end{table}
\clearpage 

%%%%%%%%%%%%%%%%%%%%%%%%%%%%%%%%%%%%%%%%%%%%%%%%%%%%%%%%%%%%%%%%%%%%%%%%%%%%%%%
\subsection{Ground and Flight Safety Risks}

\begin{table}[htbp]
\centering
\caption{Ground and flight safety risks}
\label{tab:safety_risks}
\begin{tabular}{p{4cm} p{5cm} p{4cm}}
\toprule
\textbf{Risk} & \textbf{Key Characteristics} & \textbf{Mitigation} \\
\midrule
Use of a pressurised vessel with 3 bar difference beween outside and inside. & Material: Aluminium & Pressure test will be conducted up to 4 bar in line with the Esrange safety manual. \\
\bottomrule
\end{tabular}
\end{table}
\clearpage

%%%%%%%%%%%%%%%%%%%%%%%%%%%%%%%%%%%%%%%%%%%%%%%%%%%%%%%%%%%%%%%%%%%%%%%%%%%%%%%
\subsection{Electrical Interfaces}

\begin{table}[htbp]
\centering
\caption{Electrical interfaces applicable to BEXUS}
\label{tab:electrical_interfaces_bexus}
\begin{tabular}{p{10cm} p{5cm}}
\toprule
\multicolumn{2}{l}{\textbf{BEXUS Electrical Interfaces}} \\
\midrule
E-Link Interface: E-Link required? Yes/No & Yes\\
Number of E-Link interfaces: & 1 \\
Number of required IP addresses: & 2 (one for the experiment and one for the ground station)\\
Data rate -- downlink (max. and average): & max: 100 Kbit/s  (average: 33 Kbit/s) \\
Data rate -- uplink (max. and average): & max: 10 Kbit/s \\
Interface type (RS-232, Ethernet): & Ethernet \\
\midrule
Power system: Gondola power required? Yes/No & Yes\\
Peak power and current consumption: & 65.05 W and 2.26 A \\
Average power and current consumption: & 23.11 W and 0.8 A \\
Total power and current consumption after lift-off & 138.67 Wh and 8.59 Ah \\
\midrule
Power system: Experiment includes batteries? Yes/No & No\\
Type of batteries: & \\
Number of batteries (and S x P): & \\
Capacity (1 battery): & Ah \\
Voltage (1 battery): & V \\
\bottomrule
\end{tabular}
\end{table}
\clearpage

%%%%%%%%%%%%%%%%%%%%%%%%%%%%%%%%%%%%%%%%%%%%%%%%%%%%%%%%%%%%%%%%%%%%%%%%%%%%%%%
\subsection{Launch Site Requirements}
This is a very preliminary assessment and might change in later versions of the SED.
\begin{table}[htbp]
\centering
\caption{Launch Site Requirements}
\label{tab:launch_site_req}
\begin{tabular}{lll}
\toprule
\textbf{Category} & \textbf{Item} & \textbf{Amount} \\
\midrule
Workstation&Table & 3 \\
 &Chair & 10 \\
 &Power Supply 0-30V & 3\\
 &Power Extension Cord & 2 \\

\bottomrule
\end{tabular}
\end{table}
\clearpage

%%%%%%%%%%%%%%%%%%%%%%%%%%%%%%%%%%%%%%%%%%%%%%%%%%%%%%%%%%%%%%%%%%%%%%%%%%%%%%%
\subsection{Flight Requirements}
\begin{table}[htbp]
\centering
\caption{Flight Requirements applicable to BEXUS}
\label{tab:flight_req_bexus}
\begin{tabular}{p{5cm} p{10cm}}
\toprule
\textbf{Optimal altitude} & 22km \\
\textbf{Preferred path} & ascent \\
\textbf{Minimum float time} & not important\\
\textbf{Ground track length} & not important\\
\textbf{Light/dark condition} & no preference but should be uniform during ascent if possible\\
\bottomrule
\end{tabular}
\end{table}
\clearpage

%%%%%%%%%%%%%%%%%%%%%%%%%%%%%%%%%%%%%%%%%%%%%%%%%%%%%%%%%%%%%%%%%%%%%%%%%%%%%%%
\subsection{Accommodation Requirements}

The experiment needs access to fresh air and the gas in- and outlets should not be covered by the canvas. They should be directed away from the vehicle and other experiments to minimize interference and contamination of the intaken air. Therefore the experiment should be accommodated in a corner of the gondola.

%%%%%%%%%%%%%%%%%%%%%%%%%%%%%%%%%%%%%%%%%%%%%%%%%%%%%%%%%%%%%%%%%%%%%%%%%%%%%%%
\section{Preparation and Test Activity at Esrange}
\label{sec:ESRANGE_activity}

This will be filled out in later versions of the SED.


%\begin{table}[htbp]
%\centering
%\caption{List of planned activities}
%\label{tab:planned_activities}
%\begin{tabular}{lllll}
%\toprule
%\textbf{Time/Day} & \textbf{Main Task} & \textbf{Description} & \textbf{Responsible} & \textbf{Duration [h:m]} \\
%\midrule
%& & & & \\
%\bottomrule
%\end{tabular}
%\end{table}
\clearpage


%%%%%%%%%%%%%%%%%%%%%%%%%%%%%%%%%%%%%%%%%%%%%%%%%%%%%%%%%%%%%%%%%%%%%%%%%%%%%%%
\section{Timeline for Countdown and Flight}
\label{sec:timeline}

This will be filled out in later versions of the SED.

%\begin{table}[htbp]
%\centering
%\caption{Example of timeline of the experiment events}
%\label{tab:timeline}
%\begin{tabular}{lll}
%\toprule
%\textbf{Time [s]} & \textbf{EVENT} & \textbf{Function} \\
%\midrule
%T - .. s & POWER ON & \\
%& & \\
%& & \\
%& & \\
%\bottomrule
%\end{tabular}
%\end{table}
\clearpage



\section{Post-Flight Activities}

This will be filled out in later versions of the SED.

%Celebrate!

%%%%%%%%%%%%%%%%%%%%%%%%%%%%%%%%%%%%%%%%%%%%%%%%%%%%%%%%%%%%%%%%%%%%%%%%%%%%%%%\newpage

\chapter{DATA ANALYSIS AND RESULTS}
\label{sec:data_analysis_results}

\section{Data Analysis Plan}

The focus of attention for \gls{mirage} is to demonstrate the applicability of the \gls{ndir} sensor for low-cost and low-complexity methane measurements in the stratosphere.
Since the biggest trade-off is resolution in this endeavour, the uncertainty of the measurement will be assessed with great care.
A lot of this will be done before the flight by carefully calibrating the sensor in the temperature and pressure range of the measurement chamber for different methane concentrations and humidity levels.
If possible, a machine learning algorithm will be trained to reduce the impact of noise on the measurements, similar to \cite{gia_accurate_2025}. This is an important prerequisite for the quality of the data collected during the flight.

To further demonstrate the reliability of the measurement data, the focus of post-flight data analysis is on comparison of the vertical methane profile from ascent and descent with other sources of methane data.
For this it is intended to use:

\begin{enumerate}
    \item Radiative transfer models (like RTTOV)
    \item Chemistry and Transport models (like MOCAGE)
    \item Satellite data
    \item In-situ measurement data from other groups
\end{enumerate}

\gls{rtm} like RTTOV can provide reference concentration profiles of atmospheric gases like methane and have been compared to \gls{tdlas} measurements before.
\gls{ctm} like MOCAGE assimilate satellite data into their calculations and give an additional gauge of what should be expected for the data collected by MIRAGE, although they have shown worse correspondence with in-situ measurements \cite{joly2020}.
The use of satellite data in addition to their indirect use via \gls{ctm} needs to be assessed.
This and the familiarisation with \gls{rtm} and \gls{ctm} will be part of a bachelor thesis project within \gls{mirage}, that is planned for several team members.

In-situ measurements were conducted with \gls{tdlas} from ESRANGE in 2016 and provide a direct comparison with more expensive but also more accurate laser-based methods \cite{Khair2017}.
The comparison is of course limited because of the different times at which measurements were done.
It is still expected to get some interesting insights by including these sources into the data analysis.
Another in-situ measurement campaign conducted over France with AirCore in 2018 compared their data to \gls{ftir} measurements \cite{RINGO}.
%This is an interesting study for us as well, especially if ESRANGE can support us with LIDAR methane measurements.
%Whether this is possible is not clear yet but will be asked for.

%As a secondary goal we want to investigate, whether a correlation between %different landscapes and methane levels in the stratosphere can be drawn.
%Depending on the quality of our data and number of data points we collected %during the floating phase of the flight, we will study fluctuations and %compare them with the GPS data collected by the balloon.
%Since freshwater bodies, wetlands and industrial areas have higher methane %emissions than for example forests, we are interested whether methane %fluctuations in the stratosphere can be related to those different landscapes.
%Machine learning might be employed to help recognising those fluctuations.

%%%%%%%%%%%%%%%%%%%%%%%%%%%%%%%%%%%%%%%%%%%%%%%%%%%%%%%%%%%%%%%%%%%%%%%%%%%%%%%
\section{Launch Campaign}

This will be filled out in later versions of the \gls{sed}.

%%%%%%%%%%%%%%%%%%%%%%%%%%%%%%%%%%%%%%%%%%%%%%%%%%%%%%%%%%%%%%%%%%%%%%%%%%%%%%%
\section{Results}

It is expected that less data points with a quality lower than other in-situ missions will be received by MIRAGE. However, if the calibration was done with the intended thoroughness, the results should be comparable and correspond well to \gls{rtm}. In that case, MIRAGE would have demonstrated that NDIR technology can measure methane levels in the stratosphere and hopefully inspire the further development and application of \gls{ndir} sensors for atmospheric monitoring.

%There is a small chance that fluctuations of methane in the stratosphere can %be connected to the emissions at ground level even if very precise %measurement techniques are given. Especially, if we fly over different large %areas like wetlands or lakes and forests, we might be able to detect changes %in concentrations. Even if this goal fails, we will have collected methane %concentration data with higher accuracy at about 25-30 km height due to the %longer measurement time.

%%%%%%%%%%%%%%%%%%%%%%%%%%%%%%%%%%%%%%%%%%%%%%%%%%%%%%%%%%%%%%%%%%%%%%%%%%%%%%%
\section{Input to Campaign Report}

This section will be filled after the Campaign.

%%%%%%%%%%%%%%%%%%%%%%%%%%%%%%%%%%%%%%%%%%%%%%%%%%%%%%%%%%%%%%%%%%%%%%%%%%%%%%%
\section{Lessons Learned}

As a fresh and still growing team there were some lessons learned already. 
\begin{enumerate}
  \item It took too much time recruit new team members and key team members are still missing; for example in software design.
  \item Wrong assumptions regarding the accuracy of the K96 sensor were made early one. This should have been prevented by asking more questions earlier on and closer communication with our contacts from Senseair AB.
  \item Holidays and more generous time estimates need to be taken into consideration to prevent finishing crucial project steps like the \gls{sed} under pressure, which can lead to worse quality.
  \item It it important to keep track of progress and check in with each sub team or team members working on specific tasks to prevent delays.
\end{enumerate}
Of course there are many more lessons of more practical nature. Everyone knows a little more about pumps and their limitations now, everyone learned something about methane sensing and many more things like that. However, the biggest lesson is probably that the documentation for such a project takes a lot of time and thought but also makes you aware of many aspects that can be easily overlooked.

%%%%%%%%%%%%%%%%%%%%%%%%%%%%%%%%%%%%%%%%%%%%%%%%%%%%%%%%%%%%%%%%%%%%%%%%%%%%%%%\newpage

\chapter{ABBREVIATIONS AND REFERENCES}
\label{sec:abbreviations_references}
\section{Abbreviations}
%\glsaddall %treats all acronyms as used
\printglossary[type=\acronymtype, style=compactacronyms, title={}, toctitle={}]
 %prints used acronyms

%use \gls{asap} to get accronym, first acronym will show it spelled out, after only prints acronym with link to this list

%%%%%%%%%%%%%%%%%%%%%%%%%%%%%%%%%%%%%%%%%%%%%%%%%%%%%%%%%%%%%%%%%%%%%%%%%%%%%%%
\section{References}
\nocite{*}
\printbibliography[heading=none]



%References have to be cited to appear, this can be done with \cite{example_book}. 
%If you want references to appear without citing it first use \nocite{*} 
%above \printbibliography like above\newpage
%%%%%%%%%%%%%%%%%%%%%%%%%%%%%%%%%%%%%%%%%%%%%%%%%%%%%%%%%%%%%%%%%%%%%%%%%%%%%%%

%appendix
\appendix

\chapter{EXPERIMENT REVIEWS}
\label{app:experiment_reviews}
%\section{Preliminary Design Review - PDR}
%\label{app:exprev_PDR}

%\section{Critical Design Review - CDR}
%\label{app:exprev_CDR}

%\section{Integration Process Review - IPR}
%\label{app:exprev_IPR}

%\section{Experiment Acceptance Review - EAR}
%\label{app:exprev_EAR}\newpage

\chapter{OUTREACH AND MEDIA COVERAGE}
\label{app:outreach}

\textbf{Website} \\
Website Coming soon

\textbf{Instagram:}
\url{https://www.instagram.com/riseupp_uu/?igsh=MXhsM3lsZmMyMjh4ZA%3D%3D#}

\textbf{LinkedIn:}
\url{https://www.linkedin.com/company/riseupp-uu/posts/?feedView=all}


\textbf{Logo:} \\
The RiseUpp logo was designed in July 2025. The \gls{mirage} logo is under development.

\begin{figure}[H]
    \centering
    
\includegraphics[width=0.5\textwidth]{0_images/figures/RiseUpp_logo.png}
    \caption{RiseUpp Logo}
    \label{fig:riseupp_logo}
\end{figure}

The RiseUpp association is an official member of Astronomisk Ungdom association and a website will be created under their domain.

\gls{mirage} patches and merchandise, such as shirts, hoodies and caps will be designed in the future.

\gls{mirage} aims to appear in local newsoutlets, such as Techna (student magazine), Uppsala Nya Tidning (local newspaper) and Student radio (local radiostation)

The \gls{mirage} team will be creating a project movie documenting the progress of the project.\newpage

\chapter{ADDITIONAL TECHNICAL INFORMATION}
\label{app:add_tech_info}

\begin{figure}[htbp]
    \centering
    \includegraphics[angle=90, width=0.60\textheight]{0_images/figures/tech_drawing.png}
    \caption{Technical drawing of the system as a whole.}
    \label{fig:cad_system_overview}
\end{figure}

\begin{figure}[htbp]
    \centering
    \includegraphics[angle=90, width=0.60\textheight]{0_images/figures/tech_drawing_2.png}
    \caption{Technical drawing of the pressurized chamber.}
    \label{fig:cad_pressure_chamber_1}
\end{figure}



\begin{figure}[htbp]
    \centering
    \includegraphics[angle=90, width=0.60\textheight]{0_images/figures/tech_drawing_4.png}
    \caption{Technical drawing of the pressurized chamber.}
    \label{fig:cad_pressure_chamber_2}
\end{figure}

\begin{figure}[htbp]
    \centering
    \includegraphics[angle=90, width=0.60\textheight]{0_images/figures/tech_drawing_3.png}
    \caption{Technical drawing of experiment enclosure. Drawings are representative and the thickness of both “Main Body” and “Main Lid” will be reduced to 2 mm by buying COTS parts.}
    \label{fig:cad_enclosure_1}
\end{figure}

\begin{figure}[htbp]
    \centering
    \includegraphics[angle=90, width=0.60\textheight]{0_images/figures/tech_drawing_5.png}
    \caption{Technical drawing of experiment enclosure. Drawings are representative and the thickness of both “Main Body” and “Main Lid” will be reduced to 2 mm by buying COTS parts.}
    \label{fig:cad_enclosure_2}
\end{figure}

\begin{figure}[htbp]
    \centering
    \includegraphics[angle=90, width=0.60\textheight]{0_images/figures/tech_drawing_6.png}
    \caption{Drawing of the internal layout.}
    \label{fig:cad_internal_layout}
\end{figure}\newpage

\chapter{CHECKLISTS}
\label{app:checklists}
\newpage

\chapter{PLANNING}\newpage

\chapter{MATERIAL SAFETY DATA SHEETS}\newpage 

\chapter{THERMAL ANALYSIS BACKGROUND}
\label{app:thermal_background}
Assuming energy balance
\begin{align}
    \dot{E}_{\mathrm{system}}=\dot{E}_{\mathrm{in}}-\dot{E}_{out}\qc
\end{align}
the first law of thermodynamics 
\begin{align}
    \Delta E_\mathrm{system}=Q-W\qc
\end{align}
can be used to relate heat fluxes in and out of a system to each other. Here, $Q$ is the net heat added to the system and $W$ is the work done on the system. Assuming that no work is being done (no volume expansion), and that only the internal energy (as opposed to e.g. kinetic energy) is of interest, this becomes
\begin{align}
    \Delta {E}_\mathrm{system}&=Q\\ 
    &=m_s\cdot c_{p,s} \cdot \Delta {T}_s\qc
\end{align}
where the subscript $_s$ signals system parameters, $m$ is the mass, $c_p$ the specific heat capacity and $T$ the temperature. Finally, this gives 
\begin{align}
    m_s\cdot c_{p,s} \cdot \dot{T}_s &= \dot{E}_{\mathrm{in}}-\dot{E}_{out}\\ 
    &= \dot{Q}_\mathrm{in}-\dot{Q}_\mathrm{out}
\end{align}
For simplicity, in the following the $\dot{}$ is dropped from $Q$. 

Potential contributors to the heat flux into the system $Q_\mathrm{in}$ are 
\begin{itemize}
    \item heat generated by electronics 
    \item heat generated by pumps 
    \item heat generated by pressurisation of air 
    \item heat through absorbance of solar radiation,
\end{itemize}
while contributors to heat flux out of the system $Q_\mathrm{out}$ are 
\begin{itemize}
    \item convective cooling via ambient air
    \item radiative cooling via emittance of infrared radiation.
\end{itemize}
In the following, a derivation of the individual contributions will follow. 

\paragraph{The heat generated by compressed air} must be analysed in two stages, as there are two compression stages in the system. The first is facilitated by the vacuum pumps, pressurising ambient air to 1 bar, the second by the compressor, pressurising to 3 bar. The discharge temperature is defined as 
\begin{align}
    T_\mathrm{out}&=T_\mathrm{in}\qty(\frac{p_\mathrm{out}}{p_\mathrm{in}})^{\frac{n-1}{n}}\qc
\end{align}
where $T_\mathrm{out}$ is the discharge temperature (air temperature leaving the pumps), $T_\mathrm{in}$ is the temperature of the air on the suction side, $p_\mathrm{in}$ is the air pressure on the suction side and $p_\mathrm{out}$ is the discharge pressure. $n$ is the polytropic factor, determined by the pump and gas. If the process were isentropic, the factor would be equal to the isentropic factor of air $\gamma=1.4$. For the diaphragm pump, due to negligible friction, $n<\gamma$ as the air is dumping part of its heat into the pump.
Conversely, the compressor, as a piston pump, might generate significant heat, further heating the air, such that $n>\gamma$.
Although the discharge temperature can reach up to 500°C (assuming $n=\gamma$), the actual heat capacity of that air diminishes, as the air density decreases. Thus the total heat energy carried by the air can be calculated as  
\begin{align}
    Q_\mathrm{heat} &= \dot{m}c_p\Delta T\\ 
    &= \rho I_\mathrm{pump} c_p\Delta T\qc
\end{align}
where $\rho$ is the air density, $I_{pump}$ the volumetric pump flow rate, $c_p$ the specific heat of air, and $\Delta T$ the temperature difference between suction side and discharge side. To determine the final discharge temperature, it needs to be calculated for the diaphragm pumps and based on that for the compressor.

\paragraph{The heat generated by the pumps} themselves can be estimated from their power requirements and an efficiency factor. The diaphragm pumps will probably have a higher effiency than the piston compressor pump, since they generate less friction. The total heat is then estimated as 
\begin{align}
    Q_\mathrm{heat} &= \epsilon P_\mathrm{draw}\qc
\end{align}
where $\epsilon$ is the effiency factor ranging between 0 and 1 and $P_\mathrm{draw}$ is the electrical power consumption of the pump. Note, that the operations schedule implies $P_\mathrm{draw}$ to vary with time.

\paragraph{The heat generated by other electric components} is similarly calculated. Given in \tref{tab:power_consumption} is the total power draw from the BEXUS system, as well as the total consumption. This is already assuming an effiency factor, such that the difference between those values is assumed to be emitted as heat. Thus,
\begin{align}
     Q_\mathrm{heat} &= P_\mathrm{draw}- P_\mathrm{consumed}   
\end{align}

\paragraph{The heat absorbed through solar radiation} depends on the system material, surface, and dimensions, as well as on weather conditions like cloudiness and time/solar zenith angle and altitude. In general the total heat flux is comprised of three components: direct flux from the sun, secondary solar flux reflected from Earth, clouds, and atmosphere and the heat of the earth in the \gls{ir}. Thus,
\begin{align}
    q &= q_\mathrm{sun}+q_\mathrm{reflected}+q_\mathrm{IR\,earth}\\
    &=\alpha_s (I_s \cos{\theta}+ I_s a F_{se} ) + \alpha_\mathrm{IR}(q_\mathrm{earth} F_{se})\qc
\end{align}
where $\alpha_s$ is the solar absorptivity of the surface, $I_s$ is the solar intensity received, $\theta$ is the angle between surface normal and sun, $a$ is the albedo coefficient of the Earth and atmosphere, $F_{se}$ is the view factor relating the amount of surface viewing in direction of the Earth, $\alpha_\mathrm{IR}$ is the \gls{ir} absorptivity of the surface, and $q_\mathrm{earth}$ is the \gls{ir} heat flux emitted by the Earth. Most of these values are or can be assumed to be roughly constant with respect to height. As the chassis and chamber are made out of aluminum the following analysis assumes $\alpha_{IR}=\dos{\epsilon}=0.1$ and $\alpha_s=0.4$.
The most dramatic change with altitude is seen in $I_s$.\par 
The atmosphere is absorbing part of the solar radiation incident on the top atmospheric layers $G_{sc}$. This can be modeled via the Lambert-Beer law as 
\begin{align}
    I_s(h) = G_{sc}\tau_0^{\frac{p(h)}{p_0}\frac{1}{\cos{\phi}+  0.505072 * (96.07995 - \phi) ^ {-1.6364}
}}\qc
\end{align}
where $\tau_0$ is the transmittance of solar light at ground level, $p_0$ is the atmospheric pressure at ground level, $p(h)$ is the pressure at a given altitude, and $\phi$ is the zenith angle. This formula assumes zero clouds. Since the minimum expected zenith angle in Kiruna in October is 77.11°, this value will be used by the following analysis.\par 
The final heat load is determined by multiplying the heat load for individual surfaces and summing over them, i.e.
\begin{align}
    Q_\mathrm{heat}=\sum_i q_i A_i.
\end{align}

\paragraph{The heat lost via convective cooling} via ambient air depends on the ambient air density and on system geometry. In general, 
\begin{align}
    Q_\mathrm{heat,A} = hA(T_\mathrm{system}-T_\mathrm{ambient})\qc
\end{align}
where $A$ is the area of the respective surface, $h$ is the convection coefficient and $T_\mathrm{ambient}$ is the temperature of the ambient air. The convection coefficient depends on the air density and can be modelled as 
\begin{align}
    h = h_0 \qty(\frac{p}{p_0})^n\qc
\end{align}
where $h_0$ defines some reference convection coefficient on the ground and is usually $h_0=15\,\mathrm{W/m^2/K}$. 
$n$ is usually set to 0.5 for balloon applications with laminar flow (true during ascend and floating phase, if the wind speed is not extremely high). The total heat load is again defined by adding the individual heat loads of the surfaces together.

\paragraph{The heat lost via blackbody radiation} is defined via the Stefan-Boltzmann law 
\begin{align}
    Q_\mathrm{heat,A}&= A\varepsilon\sigma (T_\mathrm{system}^4-T_\mathrm{ambient}^4)\qc
\end{align}
where $\varepsilon$ is the emissivity and $\sigma$ is the Stefan-Boltzmann constant. The total heat load is again defined by adding the individual heat loads of the surfaces together.

\paragraph{The lumped system approach} is used to model subsystems with thermally distinct behaviours. Here, this means that the system is split thermally into two parts: 
\begin{itemize}
    \item the box/chassis containing the electronics and components, 
    \item the pressure chamber within,
    \item and potentially isolation if considered (see cold and hot case in \sref{sec: thermal cold case} and \sref{sec: thermal hot case}, respectively).
\end{itemize}
The division of the system into the lumps chassis and chamber is done as the pressure chamber is to be kept at a different temperature than the rest of the system. First, the temperature of the compartment/chassis containing all electronics and components is calculated based on a constant temperature of the pressure chamber. Then, the heat flux necessary to keep the pressure chamber at the desired temperature is determined. This means that the heat flux between chamber and chassis is added as a heat load to both. 
This transfer flux is defined as 
\begin{align}
    Q_\mathrm{transfer} = F_g (T_\mathrm{chamber}-T_\mathrm{chassis}) + \epsilon_\mathrm{chamber}\sigma A_\mathrm{chamber}\qty(T_\mathrm{chamber}^4-T_\mathrm{chassis}^4)
\end{align}
Here, $F_g$ is a coefficient determining the efficiency of conductive and radiative heat transfer between chassis and chamber. While this coefficient depends on the exact geometry and material of conncecting bolts and surfaces, as well as on the air density inside the chassis, it is approximated at $F_g\approx 1.5$ W/K. The second term is responsible for modeling the radiative heat transfer between chamber and chassis. As the chassis has a much greater surface area than the chamber, it is modelled as a heat sink. This means that the efficiency of the radiative transfer is solely determined by chamber emissivity $\epsilon$ and surface area $A$.

\paragraph{Insulation} is modelled as a white styrofoam box of 5 cm thickness, completely encapsulating the chassis. The chassis is no longer subject to solar radiation, air convection, and can no longer loose heat via blackbody radiation. Thus, $Q_\mathrm{rad\,in}=Q_\mathrm{rad\,out}=Q_\mathrm{conv}=0$. Instead, the chassis is interacting with the styrofoam box via 
\begin{align}
    Q_\mathrm{insulation}=\kappa_\mathrm{styro}\frac{A_\mathrm{styro}}{d_\mathrm{styro}}\qty(T_\mathrm{styro}-T_\mathrm{chassis})\qc
\end{align} 
where $\kappa_\mathrm{styro}$ is the thermal conductivity and set to be 0.033 $\mathrm{W\,m^{-1}\,K^{-1}}$, $A_\mathrm{styro}$ is the insulation surface area, $d_\mathrm{styro}$ its thickness and $T_\mathrm{styro}$ its outer temperature. The outer temperature in turn is defined via 
\begin{align}
    Q_\mathrm{rad\,in} = Q_\mathrm{rad\,out}+Q_\mathrm{conv} + Q_\mathrm{insulation}.
\end{align}
This means that a steady-state temperature of the styrofoam is assumed. This holds true if the styrofoam's thermal capacity is small compared to that of the chassis. Since styrofoam is lightweight and the chassis is made out of aluminium this is justified.
The \gls{ir} emissivity is set to $\epsilon=\alpha_\mathrm{IR}=0.6$ and the solar absorptivity to $\alpha_s=0.15$.\par 

All in all this means that for the steady state assumptions used in the cold and hot cases, the following two equations need to be solved:
\begin{align}
    Q_\mathrm{ele}+Q_\mathrm{pump}+Q_\mathrm{transfer} + Q_\mathrm{rad\,in} + Q_\mathrm{insulation} &= Q_\mathrm{conv} + Q_\mathrm{rad\,out} \label{eq:heat_balance_compartment}\\ 
    Q_\mathrm{thermal\,system}+Q_\mathrm{compressed\,air} &=  Q_\mathrm{transfer} \label{eq:heat_balance_chamber}
\end{align} 
\eqref{eq:heat_balance_compartment} describes the heat balance of the bigger chassis and \eqref{eq:heat_balance_chamber} describes the heat balance of the pressure chamber within. Note, that depending on the assumptions in the cold and hot cases, some of the terms might be 0.%%%%%%%%%%%%%%%%%%%%%%%%%%%%%%%%%%%%%%%%%%%%%%%%%%%%%%%%%%%%%%%%%%%%%%%%%%%%%%%

\end{document}
